% Standalone research note. Intentionally supplied as uncompiled TeX.
% No external inputs, bibliography databases, figures, or generated assets are required.
\documentclass[11pt]{article}
\usepackage[T1]{fontenc}
\usepackage[utf8]{inputenc}
\usepackage[margin=1in]{geometry}
\usepackage{amsmath,amssymb,amsthm,mathtools,enumitem}
\usepackage[hidelinks]{hyperref}
\newtheorem{theorem}{Theorem}[section]
\newtheorem{proposition}[theorem]{Proposition}
\newtheorem{lemma}[theorem]{Lemma}
\newtheorem{corollary}[theorem]{Corollary}
\theoremstyle{definition}
\newtheorem{definition}[theorem]{Definition}
\newtheorem{example}[theorem]{Example}
\theoremstyle{remark}\newtheorem{remark}[theorem]{Remark}
\newcommand{\E}{\mathcal E}
\newcommand{\C}{\mathcal C}
\newcommand{\N}{\mathcal N}
\newcommand{\R}{\mathbb R}
\newcommand{\id}{\operatorname{id}}
\newcommand{\im}{\operatorname{im}}
\newcommand{\Ker}{\operatorname{Ker}}
\newcommand{\ev}{\operatorname{ev}}
\newcommand{\RegEpi}{\twoheadrightarrow}
\newcommand{\Beh}{\mathsf B}
\newcommand{\Coalg}{\mathbf{Coalg}}
\newcommand{\Obs}{\mathbf{Obs}}
\newcommand{\Poly}{\mathbf{Poly}}
\newcommand{\Tests}{\mathsf T}
\newcommand{\grad}{\operatorname{grad}}
\newcommand{\Ob}{\mathsf{Obj}}
\newcommand{\VF}{\mathsf{Fld}}
\newcommand{\Red}{\mathsf{Red}}
\newcommand{\Lear}{\mathsf{Lear}}
\newcommand{\Ext}{\operatorname{Ext}}
\newcommand{\Int}{\operatorname{Int}}
\newcommand{\Par}{\operatorname{Par}}
\title{Contextual Reduction and Wiring for Open Learners\\
\large Behavioral images, native optimization, and structural characterizations}
\author{}\date{5 November 2026}
\begin{document}\maketitle
\begin{abstract}
Can characterizations generated by learning themselves be composed?
We answer a precise general version: compare wiring reduced component states
with reducing the behavior of the complete wire. In a regular category,
a sound reduction has a canonical comparison span. Its monic leg records
which apparently compatible reduced states have an actual joint realization;
its regular-epic leg records further distinctions hidden by the exterior
interface. Exact lifting and exterior separation are independently necessary
and jointly sufficient for a strong comparison. Fixed interfaces and
covering squares supply lifting; executable separating probes supply
separation. Actual wiring functors preserving the relevant epimorphisms
give coherent comparison maps and composition followed by reduction.

The abstract statement becomes native in two substantial classes.
Deterministic polynomial machines, including target-request learners, admit
behavioral quotients on arbitrary sets. Smooth cotangent-feedback learners
admit them under regular response-coordinate and update-closure conditions;
a separate cometric criterion determines whether the actual gradient
optimizer descends. The latter is compatible with smooth objective assembly,
but gradient-flow descent alone does not imply finite-step descent.
Finally, the same realized comparison carries response incidence, formal
concepts, and Dowker topology. Exact behavioral adequacy is stronger than
concept or homotopy adequacy, and none of them reconstructs forgotten native
witnesses. General results drive the development; finite countermodels test
their hypotheses. The minimization comparison, categorical learner typing,
and gradient assembly have close established predecessors, which are
identified explicitly. No historical-priority claim is made.
\end{abstract}

\section{The question and the generalization being pursued}
The starting question is whether learning-generated characterizations can
be used compositionally, rather than only describing a chosen isolated
model. Earlier work in the programme identifies a useful common source:
actual interventions generate continuation observations; their equality
defines predictive states; their incidence supports formal concept analysis
(FCA) and Dowker constructions for topological data analysis (TDA)
\cite{CommonNote,LearningNote}. A learning update can expose a distinction
that present predictions conceal. Closing the observations under those
updates forces a richer state description. The question here is what happens
when such descriptions are wired together.

There are two routes. One can construct adequate states for the components
and then wire them, or wire the original systems and construct an adequate
state for the exterior behavior. The second route may admit fewer states
because the wire hides component distinctions. If an interface itself is
also reduced, the first route may admit combinations with no compatible
original realization. A satisfactory generalization must separate these
phenomena, say which maps actually exist, and identify when those maps live
in the intended native category of learners or optimizers.

This note treats that question across a general regular-category semantics,
a concrete categorical doctrine of deterministic learners, and a smooth
native doctrine with actual updates. The deterministic class is not restricted
to finite state spaces. The smooth class is not restricted to one architecture
or one loss. Small finite examples provide countermodels and certificates;
they are not the target definition of the framework.

\paragraph{The main conclusion.}
Under contextual soundness the comparison is, in general,
\[
 \text{all compatible reduced profiles}
 \ \longleftarrow\ \text{realizable reduced profiles}
 \ \longrightarrow\ \text{exterior behavioral states}.
\]
The first arrow is monic and the second regular epic. Strong commutation
requires both to be invertible. A fixed interface, or a proved boundary
covering condition, removes the first obstruction. A sufficient family of
external probes removes the second. Smooth descent and optimizer geometry
are additional native obligations, not consequences of these set-level
conditions. FCA and topology are transported along the realized comparison;
their adequacy conditions can be strictly weaker than behavioral adequacy.

\paragraph{Three categorical levels must remain distinct.}
A homwise contextual quotient identifies two parallel computations when all
admitted surrounding programs give the same answers. If contexts are closed
under substitution, this is already a congruence and the quotient inherits
composition. A state-space reduction instead replaces the internal carrier
of a particular open system, and asks whether its operations descend.
A native realization further requires an implementation in an independently
given grammar or category, with its smoothness, effects, ownership, resource,
and witness structure. A functor at one level is not automatically a
functor at the other two.

For example, for a typed wiring operation $w$ on implementations, a quotient
operation exists exactly when
\[
 \Ker\!\left(\prod_i q_i\right)
       \subseteq\Ker(q_{\mathrm{out}}w).
\]
If every exterior test with the other components fixed is an admitted local
test, changing one component at a time proves this inclusion. Defining the
local tests to be all substitution-closed contexts builds that argument into
the definition. It does not construct a smaller native state, discover an
operational language, or prove a realization theorem. We will give the
additional state-level construction and its native instances explicitly.

\paragraph{How this relates to actual prior results.}
The directed minimization comparison is established mathematics, not a
proposed discovery. Rosebrugh--Sabadini--Walters already give
$M(V)M(U)\to M(VU)$ in Proposition 14, organize minimization as a lax monad
in Proposition 46, and use composition followed by minimization in
Corollary 47 \cite{RSW}. Their abstract Proposition 43 also covers the
local-algebra construction from an idempotent lax monad. Spivak supplies
polynomial-coalgebra wiring and learner hom categories \cite{spivak}.
Injective change of coalgebraic behavior type reflects behavior under the
stated hypotheses of Bartels--Sokolova--de Vink \cite{bartels}.
Hanks--Fairbanks--Klawonn give compositional gradient assembly \cite{HFK}.
The purpose here is to integrate these mechanisms with the programme's
native, observation, and witness requirements, prove the exact distinctions,
and exhibit a broad smooth class in which the bridge can actually be used.
Neither a new name for a quotient nor generic categorical packaging is a
claim of novelty.

\paragraph{Scope of the source data.}
Interfaces, controls, readouts, and resource-sharing rules are supplied.
Optimizer memory is part of state when it affects updates. The observations
are exact and can range over all admitted finite continuations; a finite
sample is not substituted for that range. Reachability and physical access
are specified separately. Partial operations must retain observable failure
or enabledness, or be encoded by genuinely typed domains. We do not infer
these structures from arbitrary data or treat a semantic representative as
an executable source program without a realization argument.


\section{Typed continuation observations and their image}\label{g:sec:typed}
\subsection{Categorical assumptions}
Fix a locally small regular category $\E$. Thus $\E$ has finite limits,
every arrow factors as a regular epimorphism followed by a monomorphism,
and regular epimorphisms are stable under pullback. A regular epimorphism
is the coequalizer of its kernel pair. These properties, including image
factorization, are assumptions; they are not available in every category
of smooth or physically realizable systems.

Let $\Gamma$ be a small category of admitted typed interventions. A state
system is a functor $X:\Gamma\to\E$. An arrow $u:a\to b$ acts as
$X(u):X(a)\to X(b)$. Supply a set $D$ of readout labels, a source type
$s(d)\in\Gamma$, a response object $O_d\in\E$, and a morphism
\[
 h_d:X(s(d))\longrightarrow O_d\qquad(d\in D).
\]
For every $a\in\Gamma$, assume that the following particular product exists:
\begin{equation}\label{g:eq:profile}
 H(a)=\prod_{\substack{d\in D\\u:a\to s(d)}}O_d.
\end{equation}
Finite completeness does not supply these possibly infinite products.
One may assume all small products, or only the products in
\eqref{g:eq:profile}. The usual set-valued construction satisfies both
assumptions. The category $\Gamma$ includes only admitted operations;
nonrepeatability, type changes, and access restrictions must be represented
in its types and arrows rather than silently replaced by total resettable
operations on one set.

For $v:a\to b$, define $H(v):H(a)\to H(b)$ by selecting the coordinate
$(d,uv)$ at the target coordinate $(d,u:b\to s(d))$. Identities and
composition hold by coordinatewise equality. Define
\begin{equation}\label{g:eq:ev}
 \ev_a=(h_dX(u))_{d,u}:X(a)\longrightarrow H(a).
\end{equation}
These arrows form a natural transformation $\ev:X\to H$.

\begin{definition}
The \emph{contextual state image} is the objectwise image factorization
\begin{equation}\label{g:eq:quotient}
 X(a)\xrightarrow{q_a}Q(a)\xrightarrow{m_a}H(a),\qquad
 q_a\text{ regular epic},\quad m_a\text{ monic}.
\end{equation}
In $\mathbf{Set}$ it identifies states with the same answer to every
$h_dX(u)$. The labels of equal readouts can remain distinct in $H$.
\end{definition}

\begin{theorem}[Typed contextual reduction]\label{g:thm:reduction}
The objects $Q(a)$ carry unique transition maps $Q(v)$ such that
$q:X\to Q$ is natural. The maps $m_a$ form a monic natural transformation
$m:Q\to H$. Original readouts descend along $q$.

The internal relation $\Ker(q_a)$ is the largest family of equivalence
relations on the $X(a)$ that is preserved by every $X(v)$ and on which
each original $h_d$ is constant. If $e:X\to Y$ is an objectwise regular-epi
state representation carrying descended transitions and descended original
readouts, there is a unique objectwise regular-epi map $t:Y\to Q$ with $te=q$ and preserving
transitions and readouts.
\end{theorem}
\begin{proof}
For $v:a\to b$, the coordinates of $\ev_bX(v)$ are coordinates of $\ev_a$.
Consequently $q_bX(v)$ equalizes the kernel pair of $q_a$: composing with
$m_b$ reduces the assertion to the coordinate equalities, and $m_b$ is
monic. Since $q_a$ coequalizes its kernel pair, there is a unique
$Q(v)$ with $Q(v)q_a=q_bX(v)$. Cancel the epimorphisms $q_a$ to establish
identity and composition. Naturality of $m$ follows by the same cancellation.
The readout at $d$ is the $(d,\id)$ coordinate of $m_{s(d)}$.

The equality $\Ker(q_a)=\Ker(\ev_a)$ follows from $m_a$ being monic.
The preceding argument proves its transition stability. Conversely, any
family of relations stable under the transitions and respected by the
original readouts is respected by $h_dX(u)$ for every $d,u$, and is therefore
contained in $\Ker(\ev_a)$. This proves the maximality statement for
internal equivalence relations, without assuming that every equivalence
relation in $\E$ has a quotient.

For $e$, the descended transitions and readouts make every generated readout
factor through $e_a$. Thus $\Ker(e_a)\subseteq\Ker(q_a)$ and the
coequalizer property gives $t_a$. Naturality and preservation of readouts
follow by canceling $e_a$. Each $t_a$ is regular epic because $t_ae_a=q_a$
is regular epic; factor $t_a$ and use the extremality of $q_a$.
Uniqueness is the same cancellation.
\end{proof}

The direction of the universal map is $Y\to Q$, not $Q\to Y$.
The contextual image is terminal among state-surjective adequate
representations of the specified source. It is not a chosen reconstruction
of the original state. No finite presentation, experiment-selection method,
acquisition bound, or source program follows from the image factorization.

Stochastic systems need a further typing check. A stochastic kernel on a
set is not an arrow of $\mathbf{Set}$, and categories of stochastic
kernels do not acquire the stated regular-category properties merely by
being called semantic categories. One may work on a suitable distribution
state space with actual deterministic transitions, or prove a coalgebraic
quotient theorem with its required lumpability condition. Equality of all
trace distributions alone need not make a deterministic quotient of the
original Markov states into a Markov coalgebra. No such lift is implied
by the theorem above.

\subsection{The image reflector, explicitly}
For fixed $\Gamma,D,O_d$, let $\C$ have the above state systems as objects
and natural transformations preserving every original readout as arrows.
This category is equivalently the slice $\E^\Gamma/H$. Indeed a map to
$H$ determines its original readouts by its identity coordinates, and
naturality recovers every other coordinate. Since $\Gamma$ is small,
finite limits and regular image factorizations in $\E^\Gamma$ are
computed objectwise. The slice is regular, with its regular-epi/mono
factorizations inherited from $\E^\Gamma$. Its limits are slice limits;
they are not all ordinary limits of the underlying objects.

Let $\C_{\mathrm{obs}}$ be the full subcategory consisting of the objects
whose evaluation map to $H$ is monic. Then $R(X)=Q$ is a reflector
$R:\C\to\C_{\mathrm{obs}}$, with unit $q_X$. To see the universal
property directly, if $Y\to H$ is monic and $f:X\to Y$ preserves
readouts, then $\Ker(q_X)\subseteq\Ker(f)$, so $f$ factors uniquely
through $q_X$. Morphisms between observable systems are inclusions of their
behavior subobjects of $H$. Image choices make $R$ a functor, with all
choices uniquely isomorphic.

This is ordinary image reflection. Writing the construction categorically
makes its typing and coherence clear; it does not establish historical
novelty. If a different semantic category lacks these images or products,
that category requires a separate existence theorem.

\section{Wiring with a possibly reduced boundary}
We first work at one state type. The same construction can be made
objectwise for typed systems, provided all maps involved are system maps.
Consider the commuting port squares
\[
\begin{array}{ccc}
X&\xrightarrow{p}&B\\
{\scriptstyle q_X}\downarrow&&\downarrow{\scriptstyle b}\\
Q_X&\xrightarrow{\bar p}&\bar B
\end{array}
\qquad
\begin{array}{ccc}
Y&\xrightarrow{r}&B\\
{\scriptstyle q_Y}\downarrow&&\downarrow{\scriptstyle b}\\
Q_Y&\xrightarrow{\bar r}&\bar B .
\end{array}
\]
The arrows $q_X,q_Y,b$ are regular epimorphisms. The raw and
component-reduced wires are
\begin{equation}\label{g:eq:wire}
 P=X\times_B Y,\qquad \bar P=Q_X\times_{\bar B}Q_Y.
\end{equation}
They are full pullbacks, not reachable subsets or additional
compatibility-constrained subobjects unless explicitly stated.
There is a canonical map
\[
 \alpha:P\longrightarrow\bar P,\qquad
 \alpha=(q_X\pi_X,q_Y\pi_Y).
\]
Let the actual exterior tests and continuations on $P$ have contextual
image $q_P:P\twoheadrightarrow Q_P$ as above. Factor
\begin{equation}\label{g:eq:alpha-image}
 P\xrightarrow{e}L\xrightarrow{j}\bar P,
 \qquad e\text{ regular epic},\quad j\text{ monic}.
\end{equation}
The object $L$ is the image of genuinely compatible raw states. In
$\mathbf{Set}$ it is precisely the set of jointly realizable reduced
component profiles.

\begin{definition}[Soundness at this wiring]\label{g:def:sound}
The component reductions are sound for this wire and its exterior questions
when
\begin{equation}\label{g:eq:sound}
 \Ker(\alpha)\ \subseteq\ \Ker(q_P).
\end{equation}
Equivalently, every exterior continuation readout on $P$ is constant on
fibers of $\alpha$. This condition must be proved from a stated contextual
semantics, not inferred from agreement on isolated predictions.
\end{definition}

\begin{theorem}[The canonical comparison and its exact obstructions]
\label{g:thm:comparison}
Under \eqref{g:eq:sound}, there is a unique regular epimorphism
$\kappa:L\twoheadrightarrow Q_P$ with $\kappa e=q_P$. Thus the
canonical comparison under soundness has the shape
\begin{equation}\label{g:eq:span}
 \boxed{\quad \bar P\ \xleftarrow{\ j\ }\ L
                   \ \xrightarrow{\ \kappa\ }\ Q_P.\quad}
\end{equation}
The two arrows express different issues.
\begin{enumerate}[label=(\roman*)]
\item $j$ is invertible exactly when $\alpha$ is a regular epimorphism.
This is \emph{joint lifting}: all compatible reduced profiles have a
compatible raw realization.
\item $\kappa$ is invertible exactly when
\begin{equation}\label{g:eq:separation}
 \Ker(q_P)=\Ker(\alpha).
\end{equation}
This is \emph{exterior separation}: all realizable reduced joint profiles
are distinguished by the admitted exterior behavior.
\end{enumerate}
In particular an isomorphism $\theta:\bar P\xrightarrow{\sim}Q_P$
satisfying $\theta\alpha=q_P$ exists if and only if both joint lifting and
\eqref{g:eq:separation} hold. When it exists it is unique. When joint lifting
alone holds, the comparison is the canonical regular epimorphism
\begin{equation}\label{g:eq:fixed-comparison}
 \kappa:\bar P\twoheadrightarrow Q_P,
 \qquad \kappa\alpha=q_P.
\end{equation}
\end{theorem}
\begin{proof}
Since $j$ is monic, $\Ker(e)=\Ker(\alpha)$. The soundness inclusion
makes $q_P$ equalize the kernel pair of $e$. Its coequalizer gives the
unique $\kappa$. Factor $\kappa=mi$ with $i$ regular epic and $m$ monic.
Then $q_P=mie$ is regular epic, hence extremal epic; therefore $m$ is an
isomorphism. So $\kappa$ is regular epic.

The image factorization gives (i). For (ii), if $\kappa$ is invertible,
$q_P=\kappa e$ implies equality of the kernel pairs. Conversely, if the
kernel pairs agree, $e$ and $q_P$ are coequalizers of the same pair.
Their induced comparison is an isomorphism. If $\theta$ is an
isomorphism with $\theta\alpha=q_P$, then $\alpha$ is regular epic and
has the same kernel pair as $q_P$. The converse follows from (i) and (ii).
Uniqueness follows by canceling the regular epimorphism $\alpha$.
\end{proof}

An extension of $\kappa$ from $L$ to all of $\bar P$ may happen to exist
in a particular example. Without joint lifting it is not supplied by this
universal property, need not be unique, and assigns behavior to states not
certified as realizable. A state-level isomorphism cannot be rescued by such
an arbitrary extension.

\section{Verifiable conditions, not an assumed commutation law}
\subsection{Fixed boundaries and covering squares}
\begin{proposition}[Fixed-boundary lifting]\label{g:prop:fixed}
If $\bar B=B$ and $b=\id_B$, then $\alpha$ is a regular epimorphism.
\end{proposition}
\begin{proof}
Factor it as
\[
 X\times_B Y\longrightarrow Q_X\times_B Y
                  \longrightarrow Q_X\times_B Q_Y.
\]
Each arrow is a pullback of one of $q_X,q_Y$, so both and their composite
are regular epimorphisms.
\end{proof}

Retaining the exact boundary is a substantive observational obligation:
the port map $p$ must factor through $q_X$, and likewise for $r$.
Including the boundary readout and its future responses in the local
observation family is one way to prove this. A quotient which retains only
the present prediction need not retain that boundary.

\begin{proposition}[One-sided boundary covering suffices]
\label{g:prop:covering}
For a changing boundary, suppose the canonical map
\begin{equation}\label{g:eq:covering}
 (q_X,p):X\longrightarrow Q_X\times_{\bar B}B
\end{equation}
is a regular epimorphism. Then $\alpha$ is a regular epimorphism for every
$Y$ and every regular-epi $q_Y$ in a commuting square as above. The
symmetric condition on $Y$ also suffices.
\end{proposition}
\begin{proof}
Factor $\alpha$ as
\[
 X\times_B Y\longrightarrow Q_X\times_{\bar B}Y
                    \longrightarrow Q_X\times_{\bar B}Q_Y.
\]
The first map is the pullback of \eqref{g:eq:covering} along
$(\id,r):Q_X\times_{\bar B}Y\to Q_X\times_{\bar B}B$; the second is
a pullback of $q_Y$. Pullback stability proves the assertion.
\end{proof}

In sets, the covering condition says that every reduced state of $X$ can
be realized at \emph{every} exact boundary value compatible with its
reduced port. It is stronger than necessary for one fixed partner, but is
a concrete independent certificate for all partners. The necessary and
sufficient condition for the particular pair is
\begin{equation}\label{g:eq:set-lift}
 \bar p(\bar x)=\bar r(\bar y) \Longrightarrow\quad
 \exists x,y:\ q_Xx=\bar x, q_Yy=\bar y, p(x)=r(y).
\end{equation}
In a regular category this is the corresponding regular-logic sequent,
or equivalently the regular-epimorphism assertion for $\alpha$; it must
not be interpreted as global-point surjectivity in a category where
points do not detect epimorphisms.

\begin{theorem}[No nontrivial universally exact boundary quotient]
\label{g:thm:boundary-no-go}
Let $b:B\twoheadrightarrow\bar B$ be a regular epimorphism. Suppose the
map $X\times_BY\to Q_X\times_{\bar B}Q_Y$ is regular epic for all
commuting port squares with regular-epi component maps. Then $b$ is an
isomorphism. Conversely, if $b$ is an isomorphism, all those comparisons
are regular epic.
\end{theorem}
\begin{proof}
Take $X=Y=Q_X=Q_Y=B$, both component quotients and original ports equal
to $\id_B$, and reduced ports equal to $b$. The comparison is the
kernel-pair diagonal
\[
 \Delta_b:B\longrightarrow B\times_{\bar B}B.
\]
A diagonal is monic. If it is also regular epic it is invertible, so $b$
is monic. A map both regular epic and monic is an isomorphism.
The converse reduces, by changing along the isomorphism $b$, to
Proposition~\ref{g:prop:fixed}.
\end{proof}

Thus ``the same theorem for changed wire types'' requires a restriction on
the component class, a covering condition, a different compatibility
semantics, or explicit retention of the realizability image. The exactness
cannot hold universally merely because all three quotient maps are
surjective. This is a standard finite-limit obstruction, stated here as a
sharp constraint on the proposed reduction programme.

\subsection{Why substitution establishes soundness in the fixed-port case}
Here is an elementary sufficient condition with actual test content.
Work in sets, fix the boundary $B$, and let $q_X,q_Y$ retain their exact
ports. For each exterior continuation test $t:P\to O_t$, assume that
plugging a fixed compatible state $y$ into its $Y$ slot gives an admitted
local test on the fiber $p^{-1}(r(y))$, and likewise on the other side.
Assume the local quotient preserves these tests on their typed fibers.

If $(x,y)$ and $(x',y')$ have the same component quotients, then
$p(x)=p(x')=r(y)=r(y')$. Both $(x',y)$ and the two endpoints are genuine
states of $P$. Local contextual adequacy gives
\[
 t(x,y)=t(x',y)=t(x',y').
\]
Thus \eqref{g:eq:sound} holds. The same argument applies to any finite
wiring expression by changing one component at a time, provided each
intermediate replacement remains well typed.

This is the usual contextual-congruence argument, not a new theorem. Its
premises matter. With a changed boundary, the intermediate pair need not
be compatible. With partial experiments, equality of answers where both
are defined does not ensure equality of enabledness. If successful
execution or failure is an observable, enabledness must be retained in
the tests. An intervention alphabet restricted to what happens to be
executable at one state is not automatically a common test alphabet for
two states being compared.

Another sufficient route is to supply a wiring functor on the categories
of open systems such that it sends the component quotient morphisms to an
observation-preserving system map. Then that map is $\alpha$ on states,
and preservation of exterior readouts proves \eqref{g:eq:sound} directly.
The existence of that wiring functor and its behavior on quotient maps
must be checked from the actual forward, backward, update, and port laws.

\subsection{Probes certify that no further distinction is hidden}
\begin{proposition}[A separating-probe certificate]\label{g:prop:probes}
Assume soundness. Suppose there is a jointly monic family of arrows
$s_i:L\to Z_i$ and arrows $d_i:Q_P\to Z_i$ satisfying
\[
 s_i e=d_iq_P\qquad\text{for every }i.
\]
Then $\kappa$ is an isomorphism. Conversely, if $\kappa$ is an
isomorphism such a family exists, for example with $s=\id_L$.
\end{proposition}
\begin{proof}
Cancel $e$ to obtain $s_i=d_i\kappa$. Hence any pair identified by
$\kappa$ is identified by every $s_i$ and must be equal. Thus $\kappa$
is monic as well as regular epic. The converse uses
$d=\kappa^{-1}$.
\end{proof}

A useful noncircular sufficient version selects the $s_i$ in advance:
when $Q_X,Q_Y$ are observable behavioral images, take their component
behavior coordinates restricted to $L$. These are jointly monic. Exhibit
actual exterior protocols recovering them.
For example, a supplied bypass or component-selecting test mode can make
each component readout and future response externally available. The
protocols then establish the displayed factorization. Without those
capabilities, serial hiding, cancellation, synchronization, or an output
aggregation can prevent such probes. ``All components were minimal'' is
not a separation proof for the exterior system.

\section{Coherence for admitted fixed-interface wiring}
The comparison is not just a one-off set map. Its coherent scope can be
stated without positing a distributive law by name.

For each interface color $c$, take a test-system category
$\C_c=\E^{\Gamma_c}/H_c$ and its image reflector $R_c$ from
Section~\ref{g:sec:typed}. Supply a colored operad of wiring operations. For each
operation $w:(c_1,\ldots,c_n)\to c$, its semantics is an actual functor
\[
 W_w:\prod_i\C_{c_i}\longrightarrow\C_c
\]
with specified identity, substitution, and, where applicable, symmetry
coherences. Suppose $W_w$ preserves the componentwise regular
epimorphisms used here. Finite products and pullback wiring over retained
fixed interfaces provide the state-level preservation by regularity;
this does not dispense with checking the semantic wiring functor.

\begin{theorem}[Coherent quotient compositor]\label{g:thm:coherence}
For every $w$ and tuple $X_i$, there is a unique regular epimorphism
\begin{equation}\label{g:eq:kappa-operad}
 k_w:W_w(R_{c_1}X_1,\ldots,R_{c_n}X_n)
           \longrightarrow R_cW_w(X_1,\ldots,X_n)
\end{equation}
satisfying
\begin{equation}\label{g:eq:kappa-unit}
 k_w W_w(q_{X_1},\ldots,q_{X_n})=q_{W_w(X_1,\ldots,X_n)}.
\end{equation}
There is also a canonical isomorphism
\[
 R_cW_w(R_{c_1}X_1,\ldots,R_{c_n}X_n)
       \cong R_cW_w(X_1,\ldots,X_n).
\]
Thus re-reduction of the reduced wire agrees with reduction of the original
wire, even when $k_w$ itself is not invertible.
The comparison maps and these isomorphisms are natural in the $X_i$.
They obey the identity, substitution,
and symmetry coherences of the wiring operad. Every iterated comparison
from reducing the leaves first to reducing the fully wired system is the
same, after the specified wiring associators. Each $k_w$ is invertible
exactly when the wired reduced system is already observable.
\end{theorem}
\begin{proof}
Write $P=W_w(X_i)$ and $T=W_w(RX_i)$. The map
$\alpha=W_w(q_{X_i}):P\to T$ is regular epic by hypothesis and is a
morphism over the exterior behavior object $H_c$. Hence
$\Ker(\alpha)\subseteq\Ker(q_P)$, and the quotient property gives
$k_w$ with \eqref{g:eq:kappa-unit}. It is regular epic by
Theorem~\ref{g:thm:comparison}. The image of $T\to H_c$ is the same
subobject as the image of $P\to H_c$, since $\alpha$ is regular epic.
Therefore $k_w$ is exactly the image quotient of $T$, with its codomain
identified with $R_cP$. It is invertible precisely when $T\to H_c$ is
monic.

For naturality, compare the two required composites after precomposing
with $W_w(q_{X_i})$. Naturality of the units $q$ makes them equal, and
this map is epic. For operadic substitution, precompose the candidate
composites with the map obtained by applying the composite wiring to all
leaf units. This map is epic by the preservation hypotheses. Each
candidate then becomes the unit of the fully wired raw system, by
\eqref{g:eq:kappa-unit}; hence they agree. Identities and symmetry follow
by the same calculation. Equivalently, their common codomain is a
subobject of the exterior behavior object, so at most one
behavior-preserving comparison can have the required type.
\end{proof}

This is the familiar lax direction for image minimization. It should be
compared directly with the bicategorical construction in \cite{RSW},
not promoted as an independent new compositional principle. Calling it a
``distributive law'' would require a specified wiring monad and a
normalization monad on the same category, and checking all of their unit
and multiplication equations. The operadic coherence just proved is the
claim needed here; no unspecified distributive law is being assumed.

If a reduced interface is introduced at an intermediate cut, the
regular-epi preservation premise can fail. Theorem~\ref{g:thm:comparison}
still applies to each \emph{whole} realized wiring diagram. It does not
license iterated replacement of realized images by full reduced
pullbacks. Those pullbacks can add spurious states, so coherence of the
strong comparison cannot be inferred by simply erasing the images.
Additional joint-lifting or base-change exactness hypotheses must be
verified at the cuts. Trace or feedback also needs its own defined
semantics and preservation theorem; finite acyclic substitution does not
prove a feedback result.



\section{Independent obstructions and retained resources}
\subsection{Boundary reduction creates phantom compatibility}
Let $B=X=Y=\{0,1\}$, take both original port maps and both component
state quotients to be identities, and collapse $B$ to a singleton.
Then
\[
 P=\{(0,0),(1,1)\},\qquad
 \bar P=\{(0,0),(0,1),(1,0),(1,1)\}.
\]
The map $\alpha$ is the diagonal inclusion. Hence $L=P$ and joint
lifting fails even though each component state quotient is an
isomorphism. Choose the exterior readout $(i,i)\mapsto i$ and identity
updates. Then $Q_P$ has two states and $\kappa:L\to Q_P$ is an
isomorphism. The entire failure is in $j:L\hookrightarrow\bar P$.
It is not a failure of the behavior quotient or its native deterministic
realization. The incompatible pairs have no raw state to represent.

\subsection{Exterior hiding makes the comparison properly epic}
Let $B$ be a singleton, let $X=Y=\{0,1\}$ with identity updates and
locally visible bit readouts, so both local systems are observable and
$q_X,q_Y$ are identities. Let the wired exterior readout be
\[
 o(x,y)=x\mathbin{\oplus}y.
\]
Then $P=\bar P=L$ has four states, while $Q_P$ has two states. The
comparison $\kappa$ is XOR. It is surjective and not injective. This
remains a valid contextual congruence: with either other component fixed,
the local bit can be distinguished, but the selected exterior observation
of the assembled pair does not recover both bits. Consequently minimal
components need not give a minimal wired implementation. Adding an
actual test mode which reads either component would change the exterior
question language and remove this collapse.

\subsection{Isolated agreement without contextual closure is not sound}
Let $X=\{0,1\}$ have only a constant local readout and identity updates,
so its isolated quotient is a singleton. Take $Y$ and $B$ singleton.
Now admit a global context reading the bit of $X$. The wired behavior
quotient has two states, whereas $\bar P$ has one. The required kernel
inclusion is false and no map $\kappa$ with $\kappa\alpha=q_P$ exists.
The error is to use an isolated test family as though it were already
closed under the newly admitted context. Continuation closure under the
old update alphabet does not imply closure under every future extension
of the wiring or observation language.

\subsection{Reachability is an additional image, not a harmless convention}
Even at fixed boundaries, Proposition~\ref{g:prop:fixed} concerns the
\emph{full} pullback of the component state objects. If only states
reachable under a common control history are permitted, that reachable
joint state object may be a proper subobject of the product of the
individually reachable state objects. For example, two bits start at
zero and both toggle under the same single control. Each component has
two reachable states; the joint system reaches only $00$ and $11$.
Replacing its reachable state object by all four component pairs adds
spurious preparations. Apply the image-span theorem to the actual
reachable object rather than importing the fixed-boundary surjectivity
claim for a different domain.


\subsection{Sharing is part of the source interface}
For $m\geq2$, partition $m$ handles into blocks and give each block one
Boolean store, initially zero. Each handle alone has exactly the same
complete behavior: it can set and read an ordinary Boolean register.
Jointly, setting handle $i$ to one and reading handle $j$ returns one
exactly when $i,j$ alias the same store. Thus all $\operatorname{Bell}(m)$
sharing architectures have identical solo theories and distinct joint
behaviors. Independent tensor states describe only the singleton-block
architecture. The other architectures require an explicitly shared state
or a resource represented once with several legitimate handles. A quotient
of the independently presented components cannot infer this missing wiring
information. The same issue applies to tied parameters and common random
sources; their joint resource must be included before reduction.



\section{A general native deterministic doctrine}
\label{sec:poly}
The abstract image argument does not by itself supply learner operations.
Here they are constructed in a native category. Polynomial coalgebras give
an established language for open deterministic machines and their wiring
\cite{spivak,rutten}. This section uses arbitrary set-sized states, positions,
and directions. Finiteness is used only in the final computational
specialization. Dependent directions are handled directly by coalgebras;
we do not silently identify this presentation with a fixed-target-arrow
functor $\Gamma\to\E$. Constant-input Moore machines are an immediate
instance of that earlier template, with $\Gamma$ the free input-word monoid.

\begin{definition}[Dependent interface and machine]
A set-sized polynomial interface is
\[
 p(X)=\coprod_{o\in O}X^{I_o},
\]
where $O$ and every $I_o$ are sets. A $p$-machine $M$ is a set
$S$ with readout $o:S\to O$ and transitions
$\delta_s:I_{o(s)}\to S$. Its coalgebra is
$\beta_M(s)=(o(s),\delta_s)$. A morphism $f:M\to N$ is a function satisfying
\[
 o_Nf=o_M,\qquad f(\delta_M(s,i))=\delta_N(f(s),i).
\]
Write $\Coalg(p)$ for this category. Ordinary Moore machines have constant
$I_o=I$. Empty input fibres represent terminal public configurations.
\end{definition}
An output is visible before choosing the next input. This makes feedback
causal: a control may depend on the current readout, but the readout does
not have to solve an instantaneous fixed-point equation involving that
control. General Mealy feedback, unbounded recursion, and concurrent shared
store do not enter by changing this terminology.

A bounded-depth experiment is an adaptive decision tree, possibly with
infinitely many branches at a node: at each observed
$o$, it either returns a Boolean or chooses an admitted $i\in I_o$ and
continues with another such tree. Choices are allowed to depend on the full
observed history. All finite depths are admitted. This is not a fixed finite
test suite. Let $\Tests_p$ be the set of these labelled experiment trees and
write $\langle s,t\rangle_M\in\{0,1\}$ for execution. Distinct labels may
denote equal tests. Put
\[
 s\equiv_p s'\quad\Longleftrightarrow\quad
 \langle s,t\rangle_M=\langle s',t\rangle_M
 \text{ for every }t\in\Tests_p.
\]
For ordinary deterministic total Moore machines this is also equality of
readouts after every fixed input word. Adaptivity changes experimental
policies and costs, but does not add distinctions beyond \emph{all} words.
For dependent inputs the decision-tree formulation also records admissibility.
The trees are semantic strategies with arbitrary output-dependent branching,
not asserted to be finite programs. Pointwise output predicates can instead
be regarded as separately indexed observations; equality of all of them is
the only property used below.

\subsection{The maximal behavioral quotient}
\begin{theorem}[Behavioral congruence on arbitrary sets]\label{p:thm:quotient}
The relation $\equiv_p$ is the largest equivalence $R$ such that
\[
 sRs'\Longrightarrow o(s)=o(s')\quad\text{and}\quad
 \delta(s,i)R\delta(s',i)\quad(i\in I_{o(s)}).
\]
Consequently
\[
 \eta_M:M\twoheadrightarrow\Beh_pM,
 \qquad s\longmapsto[s]_p,
\]
is a surjective coalgebra morphism with unique quotient transitions.
The machine $\Beh_pM$ has no two behaviorally equivalent distinct
states. If $\rho:M\twoheadrightarrow N$ is any surjective coalgebra morphism,
there is a unique coalgebra morphism $N\to\Beh_pM$ whose composite with
$\rho$ is $\eta_M$.

More generally, $\Beh_p:\Coalg(p)\to\Obs(p)$ is left adjoint to the full
inclusion of behaviorally observable machines. On morphisms,
$\Beh_p(f)([s])=[f(s)]$.
\end{theorem}
\begin{proof}
Equality under all experiments gives equality of current outputs, since
immediate Boolean predicates separate $O$. If equal-output states have
successors separated by an experiment, prefix that experiment by the
corresponding input at their common output and terminate arbitrarily at
other outputs. This separates the original states. Hence $\equiv_p$ is a
congruence. Conversely every displayed congruence preserves every bounded-depth
experiment, by induction on the experiment tree. Quotient transitions are
therefore well defined.

A coalgebra morphism preserves experiments, by the same induction. Thus
$\ker\rho\subseteq\equiv_p$, giving the asserted factorization; its coalgebra
equations follow by surjectivity. The quotient is observable because any
pair of its distinct classes has an experiment that separated their
representatives. If $N$ is observable, every morphism $f:M\to N$ is constant
on $\equiv_p$, so factors uniquely through $\eta_M$. This proves the
adjunction and functoriality.
\end{proof}
The quotient is maximal as a \emph{congruence} and minimal as a number of
behavioral states among quotients. It is not a minimum mechanism, a minimum
resource implementation, or a reconstruction of the hidden state. With an
initial state fixed, one can first restrict to reachable states; the
uninitialized result above intentionally compares all supplied preparations.



\subsection{Wiring, variance, and the directed comparison}\label{p:sec:wiring}
\begin{definition}[A wiring morphism]
For $p=\sum_o\mathcal y^{I_o}$ and $q=\sum_v\mathcal y^{J_v}$, a polynomial
map $\varphi:p\to q$ consists of
\[
 \varphi_1:O\to V,
 \qquad \varphi_o^\sharp:J_{\varphi_1(o)}\to I_o.
\]
The output translation is covariant; the input translation is contravariant.
It sends a machine to the $q$-machine $\varphi_*M$ on the same set $S$ with
\[
 o^{\varphi}(s)=\varphi_1(o(s)),\qquad
 \delta^{\varphi}(s,j)=\delta(s,\varphi_{o(s)}^\sharp(j)).
\]
On coalgebra morphisms $\varphi_*$ keeps the underlying function.
\end{definition}
For $\psi:q\to r$, the input map of the composite is
$(\psi\varphi)_o^\sharp
=\varphi_o^\sharp\psi_{\varphi_1(o)}^\sharp$; consequently
$(\psi\varphi)_*=\psi_*\varphi_*$. Thus $p\mapsto\Coalg(p)$ is a functor
$\Poly\to\mathbf{Cat}$.

Every $q$-experiment pulls back to a $p$-experiment:
\[
 \varphi^*:\Tests_q\longrightarrow\Tests_p,
 \qquad
 \langle s,t\rangle_{\varphi_*M}
 =\langle s,\varphi^*t\rangle_M.                         \tag{W}
\]
At a current output $o$, report $\varphi_1(o)$ to the external experiment
and translate its chosen $j$ to $\varphi_o^\sharp(j)$. This proves (W) by
induction, and also proves
$(\psi\varphi)^*=\varphi^*\psi^*$.
Even an external fixed word may pull back to an \emph{adaptive} source
experiment, because $\varphi_o^\sharp$ can depend on the current output.
There is generally no fixed-word-to-fixed-word translation.

\begin{theorem}[Reduction--wiring comparison]\label{p:thm:comparison}
For every $M$ and every $\varphi:p\to q$, there is a canonical
surjective $q$-coalgebra morphism
\[
 c_{\varphi,M}:\varphi_*(\Beh_pM)
       \twoheadrightarrow\Beh_q(\varphi_*M),
 \qquad [s]_p\longmapsto[s]_{\varphi}.                    \tag{C}
\]
It is natural in $M$ and satisfies
\[
 c_{\id,M}=\id,\qquad
 c_{\psi\varphi,M}
 =c_{\psi,\varphi_*M}\circ\psi_*(c_{\varphi,M}).          \tag{Coh}
\]
The following are equivalent:
\begin{enumerate}[label=\textup{(\roman*)}]
\item $c_{\varphi,M}$ is an isomorphism;
\item exterior equivalence implies source equivalence on $M$;
\item the pulled-back tests $\varphi^*\Tests_q$ separate all distinct states
of $\Beh_pM$.
\end{enumerate}
For arbitrary $\varphi$, minimization once more always gives a canonical
isomorphism
\[
 \Beh_q\bigl(\varphi_*(\Beh_pM)\bigr)
       \cong\Beh_q(\varphi_*M).                          \tag{M}
\]
\end{theorem}
\begin{proof}
Equation (W) shows $\equiv_p\subseteq\equiv_{\varphi}$. Thus (C) is
well defined and surjective. Its output and transition equations follow
from the quotient definitions. Naturality and (Coh) can be checked after
precomposition with the surjection from $S$; both sides then send $s$ to
the same final equivalence class. A bijective coalgebra morphism has a
coalgebra inverse, so the equivalences follow directly from its kernel.
Lastly any surjective coalgebra morphism preserves and reflects equality
of full behaviors: both source states and their images answer exactly
the same tests. Its induced map on behavioral quotients is therefore
bijective. Apply this to (C), and use idempotence of $\Beh_q$.
\end{proof}
The direction in (C) matters. There is not in general a map selecting a
locally minimal state from a global behavioral state, and no canonical
inverse realization. One may call (C) the lax comparison with the explicit
convention $\varphi_*\Beh_p\Rightarrow\Beh_q\varphi_*$; the displayed
orientation, not the word ``lax'', specifies the claim.

Define reduced wiring on observable machines by
$\overline{\varphi}_*M=\Beh_q(\varphi_*M)$. Formula (M) supplies coherent
composition isomorphisms
\[
 \overline{\psi}_*\overline{\varphi}_*M
       \cong\overline{\psi\varphi}_*M.
\]
They send an iterated quotient class to the class of the same original
state. Hence all associativity and unit diagrams commute. The categories
$\Obs(p)$ therefore form a pseudofunctor under reduced wiring. It is
\emph{reduced} wiring, rather than the unreduced state-level operation,
that composes by these isomorphisms without extra hypotheses.

\subsubsection{An explicit uniform losslessness criterion}
For explicitly given finite tables, the test-separation criterion can be
decided by partition refinement. The next theorem removes the machine from the
condition entirely.

\begin{theorem}[Universal behavioral reflection]\label{p:thm:uniform}
Let $p=\sum_{o\in O}\mathcal y^{I_o}$ have at least two output positions.
For a polynomial map $\varphi:p\to q$, the following are equivalent:
\begin{enumerate}[label=\textup{(\roman*)}]
\item $c_{\varphi,M}$ is invertible for every $p$-machine $M$;
\item $\varphi_1$ is injective and every
$\varphi_o^\sharp:J_{\varphi_1(o)}\to I_o$ is surjective.
\end{enumerate}
When these conditions fail, a machine with at most three states witnesses
the failure. The same criterion follows even if (i) quantifies only over finite-state
machines. When the interfaces are finite and explicit, it is decidable by
finite interface inspection, without enumerating machines.
\end{theorem}
\begin{proof}
Assume (ii), and let $R$ be exterior behavioral equivalence. For $sRt$,
output injectivity gives $o(s)=o(t)=o$. Given $i\in I_o$, choose $j$ with
$\varphi_o^\sharp(j)=i$. Exterior congruence then gives
$\delta(s,i)R\delta(t,i)$. Thus $R$ is a source congruence and is contained
in source behavioral equivalence, proving (i).

If output injectivity fails, take two states with different outputs that
have the same exterior image, each self-looping on all its admitted
inputs. They are exterior equivalent but not source equivalent.
Suppose next that output injectivity holds but some $i_0\in I_o$ is missed.
Choose $o'\ne o$, and use states $s,t,z$ with outputs $o,o,o'$.
Every transition from $s$ returns $s$. Every transition from $t$ returns
$t$, except $i_0$, which returns $z$. Every transition from $z$ returns
$z$. States $s,t$ are separated in the source by $i_0$ and a current
output test, but their exterior transitions always return to themselves,
so they are exterior equivalent. This proves necessity and the witness
bound.
\end{proof}
If $|O|=1$, every nonempty source machine has the same single behavior,
including when its sole input fibre is empty. Hence input surjectivity is
not necessary in that degenerate case. If $O=\varnothing$, there are no
nonempty source machines. These exceptions are not optional side conditions.
The positive direction is an instance of an established theorem:
componentwise injective natural transformations reflect coalgebraic
behavior \cite[Lemma 4.6]{bonchi}; the earlier weak-pullback formulation
is \cite[Theorem 3.9]{bartels}. For polynomial maps, componentwise
injectivity is equivalent to output injectivity and fibrewise input
surjectivity. Indeed evaluate at a one-element set to test outputs and
at a two-element set to test whether precomposition distinguishes every
function on $I_o$. Conversely these two conditions make every component
$\varphi_X$ injective. The tiny finite countermodels above specialize the
converse and explain its one-output exception; they are not presented as
a historical novelty claim.



\subsection{Independent networks and open learners}\label{p:sec:networks}
For interfaces $p=\sum_o\mathcal y^{I_o}$ and
$q=\sum_v\mathcal y^{J_v}$ set
\[
 p\otimes q=\sum_{(o,v)}\mathcal y^{I_o\times J_v}.
\]
Independent parallel machines have state $S\times T$, readout
$(o(s),v(t))$, and update
$(s,t),(i,j)\mapsto(\delta(s,i),\epsilon(t,j))$.
This product is the established polynomial-coalgebra combination, not an
assumption that arbitrary interacting stores are independent.

\begin{proposition}[Product and network reduction]\label{p:prop:product}
There is always a surjective coalgebra morphism
\[
 \Beh_pM\otimes\Beh_qN\twoheadrightarrow\Beh_{p\otimes q}(M\otimes N),
 \qquad([s],[t])\longmapsto[(s,t)].
\]
If every input fibre of both interfaces is nonempty, this morphism is an
isomorphism. For any finite independent network $M_1,\ldots,M_k$ and
wiring $\varphi:p_1\otimes\cdots\otimes p_k\to r$, there is a coherent
surjective comparison
\[
 \varphi_*\bigl(\Beh_{p_1}M_1\otimes\cdots\otimes\Beh_{p_k}M_k\bigr)
       \twoheadrightarrow
 \Beh_r\bigl(\varphi_*(M_1\otimes\cdots\otimes M_k)\bigr). \tag{N}
\]
In particular local reduction is always sound for these networks. Exact
commutation without further minimization requires the exterior tests to
separate the product of local behavioral states.
\end{proposition}
\begin{proof}
The product of the local congruences is a global congruence, yielding the
first comparison. Under nonempty fibres, any bounded-depth experiment separating
one component can be run in parallel with arbitrary valid controls on
the other component. Stop as soon as the first component's separating
experiment stops. Thus global equivalence implies both local
equivalences. The converse already follows from the product congruence.
This proves bijectivity. Compose the product comparison with (C) to get
(N). All coherence diagrams agree on representatives of the product
state set, and surjectivity proves their equality.
\end{proof}
Nonempty fibres are a real condition. Tensoring a machine with a terminal
machine whose input fibre is empty removes every next step from the
product. States distinguished only by future source outputs then collapse.
This is synchronized blocking, not a failure of the proof.
A stateful finite controller can itself be an independent component of
such a network. Hence the theorem includes arbitrary finite controller
memory and feedback assembled by these causal wiring maps.

\subsubsection{The target-request learner interface}
For nonempty sets $A,B$, a deterministic learner consists of a
parameter set $P$ and maps
\[
 I:P\times A\to B,\qquad
 R:P\times A\times B\to A,\qquad
 U:P\times A\times B\to P.
\]
Its polynomial interface is
\[
 L_{A,B}=B^A\times A^{A\times B}\;\mathcal y^{A\times B}.
\]
At state $p$, its readout is the pair of complete tables $(I_p,R_p)$ and
its input $(a,b)$ selects the update $U(p,a,b)$. This is the existing learner--coalgebra equivalence
\cite{spivak}. The correct local quotient must preserve future
\emph{implementation and request} tables. Preserving predictions alone
need not preserve a composite learner's backward information.

For $L:A\to B$ and $K:B\to C$, serial learner composition is a polynomial
wiring
$L_{A,B}\otimes L_{B,C}\to L_{A,C}$ with output and input translations
\begin{align*}
 I_{K\circ L}(a)&=I_K(I_L(a)),\\
 R_{K\circ L}(a,c)&=R_L(a,R_K(I_L(a),c)),\\
 (a,c)&\longmapsto
 \bigl((a,R_K(I_L(a),c)),(I_L(a),c)\bigr).
\end{align*}
Thus (N) applies directly, including the actual requests used by the
updates. It does not assert that this serial wiring is uniformly lossless:
composing an observer that hides output distinctions can create extra
global equivalences. A set quotient need not carry a smooth structure or belong to a fixed
source architecture. The next section proves a separate native smooth
criterion rather than inferring one from this set-based construction.


The table-valued readout is an extensional interface description, not a claim
that an infinite table can be physically inspected in one step. Equality of
readouts means equality of every pointwise implementation and request query.
If requests are inaccessible and no admitted context can recover them, the
chosen physical observation regime can be coarser; it must be analyzed with
that regime rather than silently called the same contextual quotient.

\subsubsection{The finite specialization supplies tests of the general theorem}
For a machine with $n\geq1$ states and explicit finite interface, start with
$R_0=\Ker(o)$ and define
\[
 R_{k+1}=\{(s,t)\in R_0:
   (\delta(s,i),\delta(t,i))\in R_k\text{ for every }i\in I_{o(s)}\}.
\]
Every strict refinement increases the number of blocks, so the chain
stabilizes by $R_{n-1}$. Its stable value is behavioral equivalence.
A pair first separated at stage $k$ has a separating experiment of depth
at most $k$: store a differing output predicate or prepend an input leading
to a previously separated successor pair. This proves finite certificates
for a supplied table, not exact identification from noisy finite samples.
A network can have $\prod_i|S_i|$ states, so no polynomial bound in the sum
of component sizes follows from compositional soundness.

Two parameterized counterfamilies test the uniform criterion. A counter on
$\mathbb Z/m$ with increment input and full-state readout is minimal;
reporting only residue modulo $d$, where $d\mid m$, leaves exactly $d$
behavioral states. Separately, a machine on $\{0,1\}^n$ with first-bit
readout and rotate/hold controls has $2^n$ distinguishable states. Hiding
rotate while retaining every present output leaves two states. Thus output
visibility and control availability are independent requirements.

For an initialized deterministic total Moore machine, a known bound $N$
on reachable states, a known finite input alphabet, exact output comparison,
and reset access permit an exhaustive finite
identification: query words through length $2N-1$, use access words through
$N-1$ and continuation rows through $N-1$, and identify equal rows.
Shortest access paths and the preceding refinement bound prove correctness.
Without the state bound, any finite transcript from a constant-output machine
is also consistent with a longer hidden counter that eventually changes
output. These access promises are separate from existence of the quotient;
CALF supplies a much broader established account \cite{calf}.


\section{Native smooth response states}
\label{sec:smooth}
The category of smooth manifolds does not inherit the general regular image
construction used above. The required quotient must be established in its
own native category. The following conditions are stated on an actual
smooth learner before constructing its reduced operations; they therefore
supply more than a declaration that a successful quotient is closed.
They also isolate the exact point at which a finite response description
can realize the full continuation theory.

\subsection{What is supplied before the construction}
All manifolds in the smooth sections are finite dimensional, Hausdorff, second
countable, and without boundary. All maps in the smooth constructions are
smooth. Ports are Euclidean
spaces; parameter and optimizer-state spaces may be arbitrary manifolds.
This choice avoids identifying an output target with a cotangent signal and
does not require a global trivialization of a manifold cotangent bundle.
Here ``native'' means native to the explicitly declared smooth learner
categories. It is not an automatic realization in a separate source syntax
or in the programme's witness-retaining contextad model.

The independent source data are a state manifold $P$, an evaluator, and an
actual update law
\begin{equation}\label{s:eq:source}
 f:P\times\R^n\longrightarrow\R^m,
 \qquad
 U:P\times\R^n\times(\R^m)^*\longrightarrow P.
\end{equation}
Optimizer memory belongs to $P$ whenever the update uses it. A step size,
schedule, or other control may be included as an additional common argument;
we suppress it. The maps are total on the stated domains. Local flows fit
only after common admissible domains are supplied, as discussed below.
The external backward signal is specifically
\begin{equation}\label{s:eq:backward}
 b_f(p,x,\alpha)=(D_xf(p,x))^*\alpha\in(\R^n)^*.
\end{equation}
This is reverse differentiation, not a target-valued request function.
An independently supplied request or backward map requires its own descent
condition and is not covered by automatic differentiation of $f$.

For fixed $a=(x,\alpha)$, write $U_a(p)=U(p,x,\alpha)$. The base observations
are the scalar coordinates $f_j(-,x)$ for every admitted input $x$. The
response family is
\begin{equation}\label{s:eq:responses}
 \mathcal H=\{f_j(-,x)\circ U_{a_k}\circ\cdots\circ U_{a_1}:
              k\geq0\}.
\end{equation}
The stated interface permits forward queries at every Euclidean input.
Equality of those functions at a state also gives equality of their input
derivatives, hence of every reverse response $b_f$ at that state. With this
full query interface, contextual equivalence is equality on all of
$\mathcal H$. Finite adaptive protocols add no distinctions: equal previous
observations cause a deterministic protocol to choose the same next
control. Randomized protocols with the same externally chosen randomness
give the same conclusion path by path. Neither reset access nor the ability
to estimate these real-valued responses from noisy data follows from this
definition. Restricting update controls restricts the words in $\mathcal H$
while all forward queries may remain available. Restricting the query-input
family as well requires adding backward responses to the base observations,
or separately proving that the remaining queries determine them. For example,
$f(p,x)=px$ with identity updates is constant at the sole query $x=0$,
whereas $b_f(p,0,\alpha)=p\alpha$ distinguishes states. No derivative
recovery from isolated input samples is assumed here.

The candidate state is chosen from these supplied responses:
\begin{equation}\label{s:eq:candidate}
 q=(h_1,\ldots,h_k):P\longrightarrow Q\subseteq\R^k,
 \qquad h_i\in\mathcal H.
\end{equation}
The input to the result includes evidence that $Q=q(P)$ is an embedded
smooth submanifold and that $q:P\to Q$ is a surjective submersion. These are
substantive regularity assumptions. A constant-rank map need not have a
globally embedded image, and an arbitrary response equivalence need not
have a manifold quotient. We do not choose a smooth structure by declaring
the desired conclusion true.

\subsection{Smooth descent and a source differential test}
\begin{lemma}[Parametric smooth descent]\label{s:lem:descent}
Let $q:P\to Q$ be a surjective submersion, $A$ a manifold, and
$H:P\times A\to N$ a smooth map. There is a unique smooth
$\bar H:Q\times A\to N$ with $H=\bar H\circ(q\times\id_A)$ if and only if
$H(p,a)=H(p',a)$ whenever $q(p)=q(p')$.
\end{lemma}
\begin{proof}
Fibre constancy and surjectivity give a unique set map. A submersion has a
smooth local section $s:V\to P$ near every point of $Q$. On $V\times A$ the
set map equals $H(s(z),a)$, so it is smooth. Necessity is immediate.
\end{proof}

\begin{theorem}[Source conditions for a smooth contextual state]
\label{s:thm:source-test}
Take the source data \eqref{s:eq:source}--\eqref{s:eq:candidate}. Suppose that
the fibres of $q$ are connected and, for every $p,x,\alpha$ and every
$v\in\ker dq_p$, the following differential identities hold:
\begin{align}
 D_pf(p,x)v&=0,\label{s:eq:vertical-forward}\\
 d(q\circ U_{x,\alpha})_p v&=0.\label{s:eq:vertical-update}
\end{align}
Then there are unique smooth maps
\[
 \bar f:Q\times\R^n\to\R^m,
 \qquad
 \bar U:Q\times\R^n\times(\R^m)^*\to Q
\]
such that
\begin{equation}\label{s:eq:descent-equations}
 f(p,x)=\bar f(q(p),x),\qquad
 q(U(p,x,\alpha))=\bar U(q(p),x,\alpha).
\end{equation}
The backward map also descends, with
\begin{equation}\label{s:eq:reverse-descent}
 b_f(p,x,\alpha)=b_{\bar f}(q(p),x,\alpha),
 \quad
 (D_pf(p,x))^*\alpha
   =(dq_p)^*(D_z\bar f(q(p),x))^*\alpha.
\end{equation}
Moreover, $q(p)=q(p')$ holds exactly when $p,p'$ have equal responses in
$\mathcal H$. Thus the supplied $Q$, rather than only a set-theoretic
quotient, is a native smooth realization of this contextual response state.
\end{theorem}
\begin{proof}
Every connected manifold is path connected. Each fibre of a submersion is
a submanifold with tangent space $\ker dq$. Equation
\eqref{s:eq:vertical-forward} says that $f(-,x)$ has zero derivative along
each fibre, hence is constant there. Equation \eqref{s:eq:vertical-update}
says the same for the $Q$-valued map $q\circ U_{x,\alpha}$. A map with zero
derivative on a connected manifold is constant, by local coordinates and
paths. Lemma~\ref{s:lem:descent}, including the input and cotangent arguments
as parameters, gives \eqref{s:eq:descent-equations} and smoothness. Differentiating
the forward identity in $x$ and in $p$ gives \eqref{s:eq:reverse-descent}.

An induction on update words proves that equal $q$-states give equal values
of every response. Conversely each coordinate $h_i$ of $q$ is itself an
admitted response, so equality of all responses implies equality of $q$.
This proves both directions without assuming that the candidate already
is a sufficient state.
\end{proof}

The hypotheses are expressed using the given source maps and a chosen
response map, before a reduced update is constructed. They are not an
algorithm for finding $q$. The quantifiers over controls and fibres may be
hard to verify. Connected fibres are needed only for this differential
test: directly verified fibre constancy suffices without connectedness.

\begin{proposition}[A concrete failed closure gives a response witness]
\label{s:prop:witness}
Suppose the coordinates of $q$ are admitted responses. If $q(p)=q(p')$ but
$q(U_a(p))\ne q(U_a(p'))$, then a single additional admitted continuation
distinguishes $p,p'$. If \eqref{s:eq:vertical-update} fails at a vertical
vector, such a pair exists arbitrarily close to that point on its fibre.
An analogous statement holds for \eqref{s:eq:vertical-forward}.
\end{proposition}
\begin{proof}
Choose a coordinate $h_i$ that differs at the updated states. The test
$h_i\circ U_a$ is an admitted response. For the differential assertion,
integrate the specified tangent vector to a short curve in the fibre.
A nonzero derivative of $q\circ U_a$ makes it nonconstant on that curve.
The forward case uses one coordinate of $f(-,x)$.
\end{proof}
For an arbitrary statistic whose coordinates are not observable, failed
closure need not certify a new observable distinction. This is why
\eqref{s:eq:candidate} is operational input rather than decorative notation.

\subsubsection{A finite sufficient certificate in polynomial presentations}
Suppose $P=\R^d$, $f,U,q$ are polynomial, and the regularity of $Q=q(P)$
has been independently established. It suffices to produce polynomial
identities
\begin{equation}\label{s:eq:poly-cert}
 f_j(p,x)=F_j(q(p),x),\qquad
 q_i(U(p,x,\alpha))=V_i(q(p),x,\alpha).
\end{equation}
For an effective exact coefficient presentation, they can be checked by
coefficient comparison. The vector $V$ maps
$Q$ into $Q$ because it equals $q(U(p,x,\alpha))$ for any representative.
Thus its restriction is the native smooth update. No fibre connectedness
is required for this certificate.

\subsection{Open learner composition preserves certified reductions}
\begin{definition}
For fixed ports $\R^n,\R^m$, let $\Lear(n,m)$ be the category whose objects
are triples $L=(P,f,U)$ as in \eqref{s:eq:source}. An arrow
$q:L\to\bar L$ is a surjective submersion on state spaces satisfying
\eqref{s:eq:descent-equations}. The backward map is always
\eqref{s:eq:backward}, so its preservation follows from the forward equation.
Arrows compose as maps of state spaces; identities are identity maps.
\end{definition}
Composite arrows remain surjective submersions and satisfy both equations,
so this is an actual category. In contrast to Fong--Spivak--Tuy\'eras's
target-valued requests, our feedback type is a cotangent space. Their
learner 2-morphisms are close prior art for the preservation equations
\cite[Section VII-C]{FST}; parametrized lenses and reverse derivative
categories are direct prior art for the cotangent formulation
\cite{CGGWZ}.

Let $L=(P,f,U)\in\Lear(n,m)$ and $M=(R,g,V)\in\Lear(m,\ell)$. Define
their serial composite on $P\times R$ by
\begin{align}
 h(p,r,x)&=g(r,f(p,x)),\label{s:eq:serial-forward}\\
 W(p,r,x,\alpha)&=\bigl(
 U(p,x,b_g(r,f(p,x),\alpha)),
 V(r,f(p,x),\alpha)\bigr).\label{s:eq:serial-update}
\end{align}
Both updates use the pre-update states. Reverse differentiation gives
\begin{equation}\label{s:eq:serial-backward}
 b_h(p,r,x,\alpha)
   =b_f(p,x,b_g(r,f(p,x),\alpha)).
\end{equation}

\begin{theorem}[Serial and parallel reduction]\label{s:thm:learner-wiring}
The formulas above define a bifunctor
\[
 \Lear(m,\ell)\times\Lear(n,m)\longrightarrow\Lear(n,\ell)
\]
which sends $(r:M\to\bar M,q:L\to\bar L)$ to the state map $q\times r$.
Parallel product likewise defines a bifunctor
\[
 \Lear(n,m)\times\Lear(n',m')
       \longrightarrow\Lear(n+n',m+m').
\]
For any finite serial/parallel expression, reducing the primitives and then
evaluating the expression gives the same forward, backward, and update maps
as evaluating first and descending along the product of the primitive
reductions, with the canonical product-coordinate identifications.
\end{theorem}
\begin{proof}
The product of surjective submersions is a surjective submersion.
The forward equation for the composite follows by substituting the two
forward equations into \eqref{s:eq:serial-forward}. The input derivative of
the second forward equation gives
\[
 b_g(r_0,y,\alpha)=b_{\bar g}(r(r_0),y,\alpha).
\]
Substitute this and the first forward equation into
\eqref{s:eq:serial-update}. Applying the two update equations componentwise
gives exactly the update of the reduced composite. Equation
\eqref{s:eq:serial-backward} follows from the chain rule and is preserved for
the same reason. Products of maps preserve identities and composition, so
this operation is a bifunctor.

Parallel forward evaluation, backward feedback, and updating are defined
componentwise, using the canonical split of an output cotangent into its
two blocks. The same componentwise argument proves the parallel statement.
Induction proves the assertion for finite expressions. For three serial
learners, either parenthesization sends the final cotangent backward through
the same three derivatives and updates each original parameter using the
same old states. Thus the canonical map
$(P\times R)\times S\to P\times(R\times S)$ is a learner isomorphism.
The identity learner has singleton state, forward map $x\mapsto x$, and
the unique update; its canonical product unitors preserve the same maps.
These are the coherence identifications used in the final assertion.
\end{proof}

This theorem supplies the hom categories, composition bifunctors, and
coherence maps actually needed here. It does not assert an unexplained
hypergraph or contextad structure. Shared parameters are not automatically
independent product states. Cyclic feedback may require solving equations
or supplying time semantics. Neither operation is silently included in
\eqref{s:eq:serial-update}. Variable sharing for objective gradients is treated
separately in Section~\ref{s:sec:wiring}.

\subsubsection{Certified compositional reduction need not be compositional minimization}
\begin{proposition}[Serial hiding can create additional equivalences]
\label{s:prop:serial-minimal}
Even when both primitives are minimal for their own contextual interfaces,
their product state need not be minimal for the external interface of their
serial composite.
\end{proposition}
\begin{proof}
Take scalar states $p,r$, identity updates, $f(p,x)=p$, and $g(r,y)=r$.
Each primitive state is observable at its own output. The serial composite
always outputs $r$ and never changes either state. Thus $(p,r)$ and $(p',r)$
are contextually equivalent for every $p,p'$. Its minimal response state is
$r$, whereas the product of the primitive minimal states is $(p,r)$.
\end{proof}
Consequently Theorem~\ref{s:thm:learner-wiring} means ``wire a certified
reduction,'' not ``minimal response quotient is a strong monoidal
reflection for every wiring.'' Equality with a composite's minimal quotient
requires a further observability or separating-context theorem. For a
parallel product with independently controllable components and both output
ports exposed, component response tests do separate the product of minimal
states. For a serial network, the internal probes required for that argument
may be unavailable.


\subsection{The comparison is native when the global smooth quotient exists}
\begin{proposition}[Smooth comparison and exactness]
\label{s:prop:smooth-comparison}
Let $L$ be a wired smooth learner, let $q:L\to T$ be its certified
reduce-then-wire map from Theorem~\ref{s:thm:learner-wiring}, and let
$b:L\to B$ be a native smooth realization of its exterior contextual
quotient. Assume both state maps are surjective submersions. If
$\Ker(q)\subseteq\Ker(b)$, there is a unique surjective-submersion learner
arrow $\kappa:T\to B$ with $\kappa q=b$. It is an isomorphism of smooth
learners exactly when $\Ker(q)=\Ker(b)$.
\end{proposition}
\begin{proof}
Fiber constancy and Lemma~\ref{s:lem:descent} produce the unique smooth
map $\kappa$; it is surjective because $b$ is. Differentiating
$b=\kappa q$ shows that $d\kappa$ is surjective at every point, since
$db$ is and every target point of $q$ has a representative. The forward
and update preservation equations follow by precomposing them with $q$
and $q\times\id$, respectively, using the two learner-arrow equations,
and canceling these surjections. Hence $\kappa$ is a native learner arrow.
Kernel equality makes it bijective; a bijective submersion has a smooth
inverse, as is seen from its local sections. Its inverse preserves the
learner equations by cancellation. The converse follows from
$b=\kappa q$ when $\kappa$ is invertible.
\end{proof}
The extra hypothesis is real: a global smooth behavioral quotient need
not exist, as the figure-eight example below proves. When it does exist,
this proposition supplies the native counterpart of the regular-category
comparison rather than applying regular-category assumptions to manifolds.
For independent parallel learning with all component inputs and cotangents
available and each output block exposed, local separating response tests
can be run while selecting arbitrary controls on the other components.
The product of the exact local response quotients is therefore already
exact. Serial hiding can violate this separation, as above.



\section{When the actual optimizer descends and assembles}
\label{sec:gradient}
A native update preserved by a quotient need not be the update obtained by
running the same named optimizer on a reduced loss. This is the remaining
bridge between contextual reduction and categorical optimization. We first
state the exact all-loss geometric criterion, then prove its categorical
and wiring consequences, and finally separate continuous from discrete
learning. The result concerns a supplied optimizer, not selection of an
optimal solution.

\subsection{The exact condition for native gradient reduction}
A positive cometric on a manifold $P$ is a smooth bundle isomorphism
$K_P:T^*P\to TP$ that is symmetric positive definite. It is the inverse of a
Riemannian metric. Write $\grad_{K_P}L=K_P\,dL$. For a surjective submersion
$q:P\to Q$ define, at $p$ and $z=q(p)$,
\begin{equation}\label{s:eq:compressed-cometric}
 A_p=dq_p\,K_P(p)\,(dq_p)^*:T_z^*Q\longrightarrow T_zQ.
\end{equation}
There is no metric-free operation that turns a parameter covector into a
parameter velocity. Equation \eqref{s:eq:reverse-descent} alone therefore
does not imply forward gradient transport.

\begin{theorem}[All-loss gradient criterion]\label{s:thm:gradient}
Let $q:P\to Q$ be a surjective submersion with source cometric $K_P$.
The following are equivalent.
\begin{enumerate}[label=(\roman*)]
\item For every $\ell\in C^\infty(Q)$, the vector field
$-K_Pd(\ell\circ q)$ is $q$-projectable.
\item The map $p\mapsto A_p$ of \eqref{s:eq:compressed-cometric} is constant
on every fibre of $q$.
\item There is a unique positive smooth cometric $K_Q$ such that
\begin{equation}\label{s:eq:cometric-condition}
 dq_pK_P(p)(dq_p)^*=K_Q(q(p))
\end{equation}
and, for every smooth $\ell$,
\begin{equation}\label{s:eq:gradient-push}
 dq_p\bigl(-\grad_{K_P}(\ell\circ q)(p)\bigr)
       =-\grad_{K_Q}\ell(q(p)).
\end{equation}
\end{enumerate}
For a preassigned target cometric, all-loss gradient transport holds if and
only if that very cometric satisfies \eqref{s:eq:cometric-condition}.
\end{theorem}
\begin{proof}
The chain rule computes the left side of \eqref{s:eq:gradient-push} as
$-A_p\,d\ell_{q(p)}$. If every such field is projectable, then for two
points $p,p'$ in one fibre, $A_p\xi=A_{p'}\xi$ for every covector $\xi$ at
their common image. Every covector occurs as $d\ell_z$ for a globally smooth
function: extend a coordinate-linear function using a bump function equal
to one near $z$. Hence $A_p=A_{p'}$.

Fibre constancy defines $K_Q$. Local sections of $q$, in a local cotangent
trivialization, make $K_Q$ smooth. For nonzero $\xi\in T_z^*Q$, the
covector $(dq_p)^*\xi$ is nonzero because $dq_p$ is surjective, so
\[
 \xi(K_Q\xi)=((dq_p)^*\xi)\bigl(K_P(dq_p)^*\xi\bigr)>0.
\]
Symmetry is inherited in the same way. The chain-rule computation proves
the asserted transport and the remaining implications. Surjectivity gives
uniqueness. The argument also proves the final statement.
\end{proof}

Condition \eqref{s:eq:cometric-condition} is the cometric description of a
Riemannian submersion. It is a familiar geometric condition, not a new
notion of optimizer. The Euclidean specialization $DqDq^{\mathsf T}$ already
appears explicitly in the earlier learning manuscript \cite{LearningNote}.
Projectability of smooth dynamics has a general exact-lumpability theory
\cite{HA}. Initialization-dependent intrinsic metrics in learning have
substantial direct antecedents \cite{MPG}; the theorem above is a global
all-starting-state statement, so any necessary initialization or conserved
quantity must be retained in $Q$ or in the declared domain.

\begin{corollary}[The smaller condition for a specified learner]
\label{s:cor:accessible}
Suppose $f(p,x)=\bar f(q(p),x)$, and fix a target cometric $K_Q$.
The source and target parameter-gradient responses agree for every input
$x$ and output cotangent $\alpha$ if and only if
\begin{equation}\label{s:eq:accessible}
 (A_p-K_Q(z))\xi=0\quad
 \text{for all }\xi\in\mathcal W_z,\qquad z=q(p),
\end{equation}
where
\[
 \mathcal W_z=\operatorname{span}
   \{(D_z\bar f(z,x))^*\alpha:x\in\R^n,\ \alpha\in(\R^m)^*\}.
\]
If $\mathcal W_z=T_z^*Q$ everywhere, this is exactly
\eqref{s:eq:cometric-condition}.
\end{corollary}
\begin{proof}
Use the parameter identity in \eqref{s:eq:reverse-descent}, apply $dqK_P$,
and compare with the target response. The difference is the linear map in
\eqref{s:eq:accessible}; vanishing on the generating covectors is equivalent
to vanishing on their span.
\end{proof}
Thus demanding preservation of every smooth loss is a strong, convenient
uniform criterion. It is not a necessary assumption for every restricted
learning task.

\subsubsection{A genuine categorical gradient statement}
Let $\Red$ have objects $(P,K_P)$ and arrows the surjective submersions
satisfying \eqref{s:eq:cometric-condition}. Identity maps and composites are
arrows: for $P\xrightarrow{q}Q\xrightarrow{r}S$,
\[
 d(rq)K_Pd(rq)^*=dr\,(dqK_Pdq^*)\,dr^*=drK_Qdr^*=K_S.
\]
For the tensor, use product manifolds with block cometrics $K_P\oplus K_R$.
Products of arrows again satisfy the condition. The usual product associators, unitors, and swaps
are cometric-preserving diffeomorphisms, so $\Red$ is a symmetric monoidal
category with singleton unit. This is a restricted category of smooth
maps; arbitrary submersions with fixed metrics are not its arrows.
The tensor is not claimed to be a categorical product in this restricted
category: pairing two reduction arrows need not be a surjective submersion.

Define $\Ob$ as follows. An object is $(P,K_P,L)$ with $L:P\to\R$ smooth.
An arrow to $(Q,K_Q,\ell)$ is a $\Red$ arrow $q$ with $L=\ell\circ q$.
Define $\VF$ to have objects $(P,K_P,V)$ with a smooth vector field $V$,
and arrows the $\Red$ arrows satisfying $dq\,V=W\circ q$.
In $\Ob$ the tensor is product of states and sum of potentials; in $\VF$
it is product of states and product of fields. Both units have singleton
state and zero potential or field. All coherence maps are the canonical
product maps, which preserve those decorations.

\begin{theorem}[Gradient respects certified smooth reductions]
\label{s:thm:functor}
The assignments
\[
 (P,K_P,L)\longmapsto(P,K_P,-K_PdL),\qquad q\longmapsto q
\]
define a symmetric monoidal functor $\Ob\to\VF$.
\end{theorem}
\begin{proof}
Theorem~\ref{s:thm:gradient} makes each objective arrow a related-field
arrow. Identities and composites are unchanged as state maps. For product
states with the block cometric, differentiation of
$L(p)+M(r)$ gives the product of the two negative gradients. The unit and
coherence maps are preserved. These are precisely the functorial and
monoidal requirements.
\end{proof}
This is a covariant reduction statement on a stated class of arrows. It is
not the contravariant objective-reindexing theorem of
Hanks--Fairbanks--Klawonn, and it does not add a hypergraph structure to
$\Red$. It is one precise categorical setting in which a response
reduction, once shown to lie in the class, really carries an optimizer.
The tensor here need not be a categorical product in $\Red$: pairing two
arrows need not yield a Riemannian submersion. No Cartesian claim is used.

\subsection{Reduction and gradient assembly for smooth wiring}
\label{s:sec:wiring}
Independently supply the following data for $i=1,\ldots,N$:
\[
 q:P\to Q,\quad q_i:P_i\to Q_i,\quad
 w_i:P\to P_i,\quad \bar w_i:Q\to Q_i,
 \qquad q_iw_i=\bar w_iq.
\]
Here $q$ is a $\Red$ arrow, each $q_i$ is a surjective submersion, and all
wiring maps are smooth. Local source potentials have the verified form
$L_i=\ell_i\circ q_i$. Set
\begin{equation}\label{s:eq:wired-loss}
 F=\sum_i L_i\circ w_i,\qquad
 \bar F=\sum_i\ell_i\circ\bar w_i.
\end{equation}
The commutative squares are genuine compatibility conditions on the source
wiring and the proposed reduction. They fail if a wire reads a discarded
coordinate.

\begin{theorem}[Reduction commutes with covector and gradient assembly]
\label{s:thm:wiring}
For the supplied diagram,
\begin{align}
 F&=\bar F\circ q,\label{s:eq:wiring-loss}\\
 dF&=\sum_i(dw_i)^*dL_i
      =(dq)^*\sum_i(d\bar w_i)^*d\ell_i,\label{s:eq:wiring-covector}\\
 dq\left[-K_P\sum_i(dw_i)^*dL_i\right]
   &=-K_Q\sum_i(d\bar w_i)^*d\ell_i\circ q.
 \label{s:eq:wiring-gradient}
\end{align}
All covectors and differentials are evaluated at the corresponding points
of the displayed diagram. Therefore the field obtained by assembling the
source objective factors and then reducing is exactly the field obtained
by reducing and then assembling the target factors.
\end{theorem}
\begin{proof}
The commuting squares give
$L_iw_i=\ell_iq_iw_i=\ell_i\bar w_iq$, proving
\eqref{s:eq:wiring-loss}. Differentiation, addition, and the chain rule give
\eqref{s:eq:wiring-covector}. Multiply by $-dqK_P$ and use
$dqK_Pdq^*=K_Q$ to obtain \eqref{s:eq:wiring-gradient}, term by term.
\end{proof}

When local positive cometrics $K_{P_i}$ are also supplied, local fields can
be stored instead of covectors. Their metric-aware assembly is
\begin{equation}\label{s:eq:metric-assembly}
 \mathcal A_{w_i}(v_i)(p)
 =K_P(p)(dw_i)_p^*K_{P_i}(w_i(p))^{-1}v_i(w_i(p)).
\end{equation}
For $v_i=-K_{P_i}dL_i$ this equals the contribution in
\eqref{s:eq:wiring-gradient}. Merely adding arbitrarily chosen local tangent
vectors is not this construction. In Euclidean spaces with standard metrics
and linear wiring, \eqref{s:eq:metric-assembly} is the familiar transpose-map
assembly. Repeated use of a global coordinate adds all its covector
contributions. The local metrics are not needed for the covector proof;
the global metric compatibility of $q$ is the essential extra hypothesis.

\subsubsection{An independently constructible pullback class}
\begin{proposition}[Preserved-boundary transverse wiring]
\label{s:prop:pullback}
Let $q_i:(P_i,K_i)\to(Q_i,\bar K_i)$, $i=1,2$, be $\Red$ arrows.
Suppose boundary maps factor as
$b_i=\bar b_iq_i:P_i\to B$, and $\bar b_1,\bar b_2$ are transverse.
Give
\[
 Q=Q_1\times_BQ_2,\qquad P=P_1\times_BP_2
\]
the metrics induced from their respective product Riemannian metrics.
Then $P,Q$ are smooth manifolds, and the restricted map
$q_1\times q_2:P\to Q$ is a surjective Riemannian submersion.
Theorem~\ref{s:thm:wiring} applies to the two projections and summed local
potentials.
\end{proposition}
\begin{proof}
The target pullback is a submanifold by transversality. The product
$q_1\times q_2$ is a surjective submersion, and $P$ is its inverse image of
$Q$, so the restricted map is also a surjective submersion.
At a source point its vertical tangent is
$\ker dq_1\oplus\ker dq_2$. The product horizontal lift of a tangent pair
$(v_1,v_2)\in TQ$ lies in $TP$: the boundary derivatives of its components
are $d\bar b_1v_1=d\bar b_2v_2$. It is orthogonal to the vertical tangent
and has the same norm as $(v_1,v_2)$, since each $q_i$ is a Riemannian
submersion. These horizontal lifts exhaust the horizontal tangent of $P$,
proving the assertion for the induced metrics. The projection squares
commute, so Theorem~\ref{s:thm:wiring} applies.
\end{proof}

This gives a large per-diagram class, including nonlinear manifolds and
genuine shared-boundary constraints. It does \emph{not} prove that induced
pullback metrics supply a unital category of metric spans. For example,
gluing the identity $\R\to\R$ to itself gives the diagonal in
$\R^2$; its induced metric is $2\,dx^2$, not $dx^2$ under the canonical
unitor. A potential pulled back to this diagonal therefore has half the
original gradient in that coordinate. One must specify coherent apex
metrics or a different decoration transport before claiming a global
hypergraph functor of these reductions. Transversality and the factorization
of boundary maps are also assumptions, not consequences of behavioral
equivalence.

\subsection{Continuous projectability does not settle discrete learning}
For $q$-related smooth vector fields, uniqueness of solutions gives
\begin{equation}\label{s:eq:flow}
 q\circ\Phi_t=\bar\Phi_t\circ q
\end{equation}
where both flows are defined. This does not guarantee that a source solution
exists whenever the target one does, nor a common nonzero time interval for
all representatives in an unbounded fibre. A total flow-based learner needs
those additional existence and domain conditions.

On manifolds, a discrete gradient update also needs a specified retraction.
Let $R_P,R_Q$ be smooth retractions, defined on neighborhoods of the zero
sections of the tangent bundles. For a $\Red$ arrow $q$ write
\[
 H_p=\im\bigl(K_P(p)(dq_p)^*\bigr)\subseteq T_pP.
\]
This is the horizontal tangent space. A suitable discrete compatibility
condition is
\begin{equation}\label{s:eq:retraction}
 q(R_P(p,v))=R_Q(q(p),dq_pv),\qquad v\in H_p,
\end{equation}
on the declared common domains.

\begin{proposition}[The exact finite-step condition]
\label{s:prop:retraction}
Let $q$ be a $\Red$ arrow. Fix a nonzero step size $\eta$ and a prescribed
domain $\mathcal D\subseteq H$ contained in the domain of $R_P$, whose
image $(p,v)\mapsto(q(p),dq_pv)$ lies in the domain of $R_Q$.
Condition \eqref{s:eq:retraction} holds on $\mathcal D$ if and only if
the source and target gradient steps commute for every smooth $\ell$ at
every $p$ whose source increment
$-\eta K_P(p)(dq_p)^*d\ell_{q(p)}$ lies in $\mathcal D$.
\end{proposition}
\begin{proof}
The source increment is
$v=-\eta K_P(dq)^*d\ell$, hence belongs to $H_p$, and its pushforward is
$-\eta K_Qd\ell$ by \eqref{s:eq:cometric-condition}. Substitution gives
sufficiency. Conversely every $(p,v)\in\mathcal D$ has that form for the
fixed nonzero $\eta$: choose its unique cotangent preimage, scale by
$-1/\eta$, and realize that covector as the differential of a smooth
potential at $q(p)$. Applying the assumed step identity gives
\eqref{s:eq:retraction} for that $v$.
\end{proof}
The fixed-domain quantifier is essential. A loss-dependent sufficiently
small step bound is not being converted into a uniform neighborhood or a
global retraction identity.

\begin{corollary}[A native reduced open gradient learner]
\label{s:cor:native-gradient-learner}
Suppose $q:P\to Q$ is a surjective submersion, $f$ is constant on its
state fibres for every input, and the source update is the declared
retraction gradient step
\[
 U(p,x,\alpha)=R_P\bigl(p,-\eta K_P(p)(D_pf(p,x))^*\alpha\bigr).
\]
Assume \eqref{s:eq:cometric-condition} and \eqref{s:eq:retraction}, and assume
all displayed updates have common total domains as in \eqref{s:eq:source}.
Then $q$ is an arrow in $\Lear(n,m)$ to the native smooth gradient learner
\[
 \bar U(z,x,\alpha)
   =R_Q\bigl(z,-\eta K_Q(z)(D_z\bar f(z,x))^*\alpha\bigr),
\]
with the descended forward and reverse maps. If the coordinates of $q$ are
admitted response observables separating its image as in
\eqref{s:eq:candidate}, this is the exact contextual response quotient.
It composes with serial and parallel learner wiring by
Theorem~\ref{s:thm:learner-wiring}.
\end{corollary}
\begin{proof}
Apply Lemma~\ref{s:lem:descent} to $f$, then use the parameter chain rule in
\eqref{s:eq:reverse-descent}. The source increment is horizontal, its
pushforward is the displayed target increment, and
\eqref{s:eq:retraction} proves the update equation. Reverse feedback descends
by differentiation in the external input. The response argument is the
same finite-word induction and coordinate-separation proof as in
Theorem~\ref{s:thm:source-test}. The composition assertion now applies to the
constructed learner arrows, not to an assumed quotient update.
\end{proof}

\begin{proposition}[All basic gradient flows can descend while Euler does not]
\label{s:prop:euler-obstruction}
There are a global surjective submersion with connected fibres and a smooth
positive source cometric satisfying \eqref{s:eq:cometric-condition}, but for
which a source coordinate Euler step has no well-defined update on the
quotient, for every nonzero step size.
\end{proposition}
\begin{proof}
Take $P=\R^2$, $Q=\R$, $q(x,y)=x+y^2$, and $K_Q=1$. Put
\[
 h(x,y)=\binom{1-2y^2}{y},\qquad
 v(x,y)=\binom{-2y}{1},\qquad K_P=hh^{\mathsf T}+vv^{\mathsf T}.
\]
The matrix with columns $h,v$ has determinant one, so $K_P$ is positive
definite everywhere. Since $dq=(1,2y)$, we have $dq\,h=1$ and $dq\,v=0$.
It follows that $K_Pdq^{\mathsf T}=h$ and $dqK_Pdq^{\mathsf T}=1$.
Every basic gradient flow therefore projects by
Theorem~\ref{s:thm:gradient}. For the linear target potential $\ell(z)=z$,
the source coordinate Euler step is $(x,y)\mapsto(x,y)-\eta h(x,y)$.
Direct calculation gives
\begin{equation}\label{s:eq:euler-failure}
 q\bigl((x,y)-\eta h(x,y)\bigr)
       =q(x,y)-\eta+\eta^2y^2.
\end{equation}
The points $(z,0)$ and $(z-1,1)$ lie in one fibre but have distinct updated
quotient values for $\eta\ne0$. Thus there is no descended map at all,
not merely a mismatch with the target Euler formula $z\mapsto z-\eta$.
\end{proof}
This source is also an open learner with $f(x,y)=q(x,y)$, no external input,
and $U_\alpha(p)=p-\eta K_P(p)dq_p^*\alpha$. The parameter geometry is
compatible with every instantaneous loss, but the original prediction
quotient is not a learning state for its actual discrete protocol.

\subsubsection{A broad positive finite-step class}
Let $P=\R^d$, $Q=\R^k$, $q(p)=Ap+c$ with $A$ full row rank, and let
$K_P$ be constant positive definite. Set $K_Q=AK_PA^{\mathsf T}$.
For an arbitrary smooth evaluator $\bar f:Q\times\R^n\to\R^m$, take
$f=\bar f\circ(q\times\id)$ and the source update
\[
 U(p,x,\alpha)=p-\eta K_P(D_pf(p,x))^*\alpha.
\]
The quotient update is exactly
\begin{equation}\label{s:eq:affine-update}
 \bar U(z,x,\alpha)
   =z-\eta K_Q(D_z\bar f(z,x))^*\alpha.
\end{equation}
Both identities follow by the chain rule and affineness of $q$. Hence this
is an unrestricted smooth nonlinear evaluator class, in arbitrary finite
dimensions, with exact discrete native reduction. It is closed under the
serial and parallel constructions with product state metrics. Standard
Euclidean target gradient descent is obtained when $AK_PA^{\mathsf T}=I$;
otherwise the induced metric or preconditioner is part of the reduced
optimizer and cannot be omitted. Minimality is still a response-separation
question, not a consequence of full row rank.
This constructs reducible learners from a supplied $q$ and reduced evaluator;
it is not a procedure for finding $q$ in an independently given architecture.

More generally, any supplied surjective submersion $q:P\to Q$ can be made
a Riemannian submersion by choosing a smooth horizontal complement to
$\ker dq$, using the pullback of a chosen target metric on that complement,
and a positive metric on the vertical bundle. Such complements and metrics
exist on these manifolds. This construction establishes that the smooth
class is nonempty and broad. It does not establish compatibility with a
\emph{previously fixed} source optimizer; changing the metric changes the
learning protocol and its response equivalence.



\subsection{A smooth source need not have a manifold behavioral quotient}
The regularity premise in Section~\ref{sec:smooth} cannot be omitted.
Consider the smooth learner with state $P=\R$, singleton input, output
\[
 f(p)=(\sin p,\sin 2p),\qquad U(p,\alpha)=p.
\]
Every finite response is a repetition of the present output. Its complete
behavioral quotient is the figure-eight image of $f$. No Hausdorff manifold
can be the target of a smooth surjection $q$ with exactly these fibers.
Indeed such a $q$ is $2\pi$-periodic and factors continuously through $S^1$.
The resulting surjection from the compact circle to a Hausdorff target is
a quotient map. Since its fibers are exactly those of the figure-eight map,
the target would be homeomorphic to that image. At the crossing the image
has four local branches, so it is not a manifold, even allowing boundary.
This proves nonexistence of a native quotient in the stated manifold class,
not merely failure of one selected coordinate map to be a submersion.
A broader stratified or diffeological doctrine may be appropriate, but its
operations and optimizer must then be constructed there.

Current prediction equality can fail still earlier, before topology matters.
For $f(p)=p^3-p$, target $b$, and Euclidean half-squared-loss step
\[
 U(p,b)=p-\tfrac1{10}(f(p)-b)(3p^2-1),
\]
states $0$ and $1$ initially predict $0$, while training on $b=1$ gives
$U(0,1)=-1/10$ and $U(1,1)=6/5$, with predictions $99/1000$ and $66/125$.
Thus current prediction is not an update congruence. Complete continuation
behavior separates these states; it does not identify them and then fail
its own deterministic quotient law.


\section{Native realization and the retained contextad level}
\subsection{What must be supplied to lift the image}
Let $U:\N\to\C$ interpret native systems in the test-system semantics.
A \emph{native realization of the quotient of $X\in\N$} consists of
$\widehat Q\in\N$, a native arrow $r:X\to\widehat Q$, and an
isomorphism $\varphi:U\widehat Q\xrightarrow{\sim}R(UX)$ such that
$\varphi U(r)=q_{UX}$. This is additional data, not a consequence of
$R(UX)$ existing in $\C$.

A simple sufficient transfer criterion is that $U$ identify $\N$ with a
full subcategory of $\C$ closed under the image reductions and under the
admitted wiring operations. Then the units and the comparisons of
Theorem~\ref{g:thm:coherence} are native arrows, and their invertibility
can be tested in $\C$ because a fully faithful interpretation reflects
isomorphisms. Closure must be established from the native doctrine, as
it was by the displayed quotient formulas for deterministic machines.
Calling an arbitrary source grammar ``quotient-closed'' without proving
it only restates the desired lifting conclusion.

If $U$ is not full, even a semantic comparison between two native objects
may fail to lift. If it forgets proof-relevant noninvertible cells,
factorization witnesses, ownership, or occurrence structure, an
isomorphism of state images does not recover those data. A strong
comparison with a native bicategory requires native functors or
pseudofunctors and their actual unit/composition cells, not only this
set-level construction. These are among the matching obligations retained by the native
context/content construction and its controlled-recovery results
\cite{RelativeNote,RecoveryNote}.

\subsection{Hidden occurrences force more native states}
Here is an obstruction in the specific native doctrine of the earlier
realization notes, not an artificially imposed prohibition on quotients.
Take one module, a singleton announcement set $J$, two hidden occurrences
$H=E=\{0,1\}$, a singleton retained set $S$, and
\[
 A_1=(1/2,1/2),\qquad B_1=\begin{pmatrix}1\\1\end{pmatrix}.
\]
All public output channels are constant. A native middle realization has
a finite set $V$, nonnegative row $A$ on $V$, a deterministic occurrence
map $f:V\to\{0,1\}$, and the prescribed occurrence law
\[
 Af=(1/2,1/2).
\]
Both occurrences have positive mass, so $f$ must take both values on
positive-mass states and $|V|\ge2$. Two states realize the data.
Nevertheless every middle state has the same public continuation
behavior; its continuation state quotient has one element.

Thus the one-state behavioral image is not a native realization of the
retained occurrence specification. This does not use a general
surjectivity requirement on occurrence maps: positive prescribed
occurrence marginals force the two preimages in this particular input.
It also does not refute contextual sufficiency for the public questions.
It shows that the native specification asks an additional question which
the public behavior quotient deliberately forgets. To retain that
question, enrich the semantic observation/witness structure or solve the
native realization problem separately.

When the native doctrine is unrestricted finite deterministic machines
and all source behaviors must be represented, the number of quotient
states is minimal among adequate surjective state representations: every
such representation maps onto $Q$ by Theorem~\ref{g:thm:reduction}. This
cardinality statement is not a constrained-rank theorem for stochastic
native mechanisms, nor a cost theorem for source programs or physical
implementations.



\section{Structural characterizations of the same comparison}
\label{sec:views}
FCA and TDA enter here through the exact source map already constructed,
not as unrelated constructions appended to an optimizer. This section is
explicitly a \emph{Set-valued Boolean interpretation} of the regular-category
comparison. An arbitrary regular category does not automatically carry
powersets, Boolean predicates, or simplicial complexes. The supplied test
family determines the incidence relation, and a chosen change of that family
can change the resulting geometry. The earlier common-generation and
learning notes already establish the basic concept and Dowker transports
\cite{CommonNote,LearningNote}; their published antecedents include
\cite{gp,ayzenberg,brun}. The issue here is the additional realizability
image inside the full reduced wire.

\subsection{Keep the realizability image in the type}
Let two actual state carriers have port maps
\[
 p_X:X\longrightarrow B,\qquad p_Y:Y\longrightarrow B.
\]
Let $q_X:X\twoheadrightarrow Q_X$, $q_Y:Y\twoheadrightarrow Q_Y$,
and $q_B:B\twoheadrightarrow Q_B$ commute with specified reduced port
maps $\bar p_X,\bar p_Y$. Form
\[
 P=X\times_B Y,\qquad
 \bar P=Q_X\times_{Q_B}Q_Y,
 \qquad
 \alpha:P\longrightarrow\bar P,
 \quad (x,y)\longmapsto(q_Xx,q_Yy).
\]
Its image factorization and the external behavioral quotient have types
\begin{equation}
 P\mathrel{\mathop{\twoheadrightarrow}^{\beta}}L
 \mathrel{\mathop{\lhook\joinrel\longrightarrow}^{j}}\bar P,
 \qquad q_P:P\twoheadrightarrow Q_P.
 \label{v:eq:span}
\end{equation}
When equality of reduced joint profiles implies external contextual
equivalence, $q_P$ factors uniquely as
\begin{equation}
 q_P=\kappa\beta,\qquad \kappa:L\twoheadrightarrow Q_P.
 \label{v:eq:kappa}
\end{equation}
Thus the meaningful comparison is $\bar P\leftarrow L\to Q_P$, not
an assumed isomorphism $\bar P\cong Q_P$. Theorem~\ref{g:thm:comparison} supplies these maps and their hypotheses.
We now analyze what follows for incidence structures.

There are two different comparisons.
\begin{itemize}
\item A designated response test $e:Q_P\to\{0,1\}$ pulls back exactly
along $\kappa$ and $\beta$. Hence it defines the same test on actual
states and realizable profiles. This is the existing sufficient-state
comparison.
\item Comparing with \emph{all} of $\bar P$ additionally requires an
interpretation $\widetilde e:\bar P\to\{0,1\}$ whose restriction to
$L$ is $e\kappa$. Such extensions and their operational laws are extra
data. Their values on unrealizable profiles do not follow from $q_P$.
\end{itemize}
Any claimed execution from an ambient profile must still lift through
$\alpha$. Equalities of Boolean observations cannot manufacture such a
lift.

\subsection{Three different adequacy tests for an inclusion}
Let $L\subseteq Z$ be nonempty sets and let $E$ be a set of labelled
Boolean tests on $Z$. In the wiring application $Z=\bar P$ and $L$ is
its realizability image. Write
\[
 r(z)=\{e\in E:zRe\},\qquad C_e=\{z:zRe\}.
\]
The context on $L$ is the restriction of $R$. Let
\begin{align*}
 \Int_Z&=\left\{\bigcap_{z\in U}r(z):U\subseteq Z\right\},&
 \Int_L&=\left\{\bigcap_{l\in V}r(l):V\subseteq L\right\},\\
 \Ext_Z&=\left\{\bigcap_{e\in A}C_e:A\subseteq E\right\},&
 \Ext_L&=\left\{\bigcap_{e\in A}(C_e\cap L):A\subseteq E\right\}.
\end{align*}
Empty intersections are $E$ on the row side and the appropriate object
carrier on the column side. These are the closed-intent and closed-extent
lattices, with the usual order reversal between them. For $S=L$ or $Z$,
write
\[
 U'_S=\{e\in E:(\forall s\in U)\ sRe\},\qquad
 A'_S=\{s\in S:(\forall e\in A)\ sRe\}.
\]
The argument type determines the direction of the prime. The supported
concept core consists of the extents $A\in\Ext_S$ with both $A\ne\varnothing$
and $A'_S\ne\varnothing$, ordered by inclusion. Equivalently it retains
exactly the formal concepts with nonempty extent and intent.
Let $K_Z$ and $K_L$
be the object Dowker complexes: a nonempty finite set is a face precisely
when it has a common test in $E$.

\begin{definition}
The realized subcontext is:
\begin{enumerate}
\item \emph{row-realizing} if for every $z\in Z$ some $l\in L$ has
$r(z)=r(l)$;
\item \emph{intersection-generating} if for every $z\in Z$ some
$V_z\subseteq L$ has $r(z)=\bigcap_{l\in V_z}r(l)$;
\item \emph{row-dominating} if for every $z\in Z$ some $l\in L$ has
$r(z)\subseteq r(l)$.
\end{enumerate}
The first condition is the row-density criterion. The second permits an
empty indexing family; the third need not preserve behavioral equality.
\end{definition}

\begin{theorem}[Exact separation of behavioral and concept adequacy]
\label{v:thm:three}
For the inclusion $j:L\hookrightarrow Z$:
\begin{enumerate}
\item The induced map of behavioral row quotients $L/{\sim_E}\to
Z/{\sim_E}$ is injective, and is bijective exactly when $L$ is
row-realizing.
\item Extent restriction
\[
 \Phi:\Ext_Z\longrightarrow\Ext_L,
 \qquad A\longmapsto A\cap L
\]
is always a surjective complete-meet map. It is an order isomorphism
of full concept lattices exactly when $L$ is intersection-generating.
Equivalently, $\Int_Z=\Int_L$.
\item Suppose the equivalent conditions in (2) hold. Restriction then
induces an isomorphism of the \emph{supported} concept posets exactly when
\begin{equation}
 (\exists z\in Z:\ r(z)=E)
 \quad\Longleftrightarrow\quad
 (\exists l\in L:\ r(l)=E).
 \label{v:eq:universal-object}
\end{equation}
This is the support flag for the bottom extent, not an unlabelled
lattice condition.
\item Under (2) and (3), $L$ is row-dominating. More generally, whenever
$L$ is row-dominating, $j:K_L\to K_Z$ is a homotopy equivalence. The
comparison is an actual simplicial inclusion with a contiguous
retraction, not merely equality of Betti numbers.
\end{enumerate}
Consequently row realization implies all three adequacies. Full concept
and geometric adequacy together do not imply row realization.
\end{theorem}

\begin{proof}
Part (1) is equality of the same labelled rows on a subset of the object
carrier: no two distinct realized rows become identified, and surjectivity
says precisely that no new row occurs.

For (2), every extent on $L$ has the form $A'_L$ for a set $A\subseteq E$,
and equals $A'_Z\cap L$. Thus $\Phi$ is surjective and commutes with
arbitrary intersections, including the empty one. We always have
$\Int_L\subseteq\Int_Z$. If every ambient row is an intersection of
realized rows, flattening intersections proves the opposite inclusion.
When the intent lattices coincide, a closed extent on either carrier is
specified by exactly the same closed intent; this gives the stated
isomorphism and its inverse.

Conversely, assume $\Phi$ is injective. Let $I\in\Int_Z$, put
$A=I'_Z$, and define $J=(A\cap L)'_L$. Then $J$ is an intersection of
realized rows and hence belongs to $\Int_L\subseteq\Int_Z$.
Moreover
\[
 \Phi(J'_Z)=J'_L=A\cap L=\Phi(A).
\]
Injectivity gives $J'_Z=A$; taking primes and using closedness gives
$J=I$. Hence $\Int_Z=\Int_L$. In particular every $r(z)$ is an
intersection of realized rows.

For (3), the corresponding concepts have the same intent, so nonemptiness
of intent is already preserved. Restriction can turn a nonempty extent
into the empty set only when the latter is the bottom extent on $L$.
Under the lattice isomorphism, the corresponding extent is the bottom
on $Z$. These bottom extents are respectively
$\{l:r(l)=E\}$ and $\{z:r(z)=E\}$. Thus their nonemptiness agrees
exactly under (3), proving the assertion. This argument covers $E$ empty:
then the supported concept posets are both empty.

For (4), choose the representations in (2). A nonempty $V_z$ supplies
any $l\in V_z$ with $r(z)\subseteq r(l)$. If $V_z$ is empty, then
$r(z)=E$ and~\eqref{v:eq:universal-object} supplies a realized universal row. This proves
row domination. Now choose $\rho:Z\to L$ with
$r(z)\subseteq r(\rho z)$ and $\rho(l)=l$ for $l\in L$.
For each face $\sigma$ and each one of its common tests, every vertex of
$\sigma\cup j\rho[\sigma]$ satisfies that test. Therefore $\rho$ is
simplicial, $\rho j=\id$, and $j\rho$ is contiguous to $\id$.
The resulting affine homotopy on each simplex gives the claimed homotopy
equivalence. Choice is only used to choose representatives; no operational
realizability of $\rho$ is asserted.
\end{proof}

\begin{remark}[What the criteria do not say]
Intersection generation concerns the \emph{marked incidence table}, not
an arbitrary accidental isomorphism of abstract lattices. Row domination
is sufficient for this concrete simplicial retraction; it is not claimed
necessary for a homotopy equivalence. A row-dominating replacement can
satisfy strictly more tests, so it is not a behavior-preserving simulator.
Even a row-realizing choice of $\rho$ need not be available as a native
operation or be definable in the source language.
\end{remark}

\subsection{The comparisons retain typed operations and their domains}
Index the preceding data by interfaces $a$. A typed operation $u:a\to b$
supplies
\begin{equation}
 f_u:Z_a\to Z_b,\qquad g_u:E_b\to E_a,
 \qquad f_u(z)R_b e\ \Longleftrightarrow\ zR_a g_u(e).
 \label{v:eq:chu}
\end{equation}
Assume $f_u[L_a]\subseteq L_b$, as holds when these are actual
realizability images and the operational squares commute. Composition has
$f_{vu}=f_vf_u$ and $g_{vu}=g_ug_v$. Do not substitute two covariant maps
for this contravariant test pullback.

For either carrier system, the induced action on closed extents is
\[
 \widehat f_u(A)=(f_u[A])''.
\]
It is left adjoint to the inverse-image map on extents: that inverse image
is closed by the Chu equation, and
$(f_u[A])''\subseteq B$ is equivalent to $A\subseteq f_u^{-1}[B]$ for
closed $B$. Its intent is exactly $g_u^{-1}[A']$. On the supported concept
core its domain therefore consists of the nonempty extents $A$ for which
$g_u^{-1}[A']\ne\varnothing$. These are the concept actions and supported
domains compared below.


For a nonempty finite simplex $\sigma\subseteq Z_a$, the maximal
permitted domain of its forward map is
\begin{equation}
 D_u^Z=\{\sigma:\exists e\in E_b\ \forall z\in\sigma,
                  f_u(z)R_be\}.
 \label{v:eq:domain}
\end{equation}
Include the empty face separately. Equation~\eqref{v:eq:chu} ensures $D_u^Z\subseteq
K_{Z_a}$. Define $D_u^L$ in the same way on $L_a$.

\begin{proposition}[A comparison on every operational domain]
\label{v:prop:domains}
Assume each $L_a$ is row-dominating in $Z_a$. Then
\begin{equation}
 K_{L_a}\longleftarrow D_u^L\longrightarrow K_{L_b}
 \quad\longrightarrow\quad
 K_{Z_a}\longleftarrow D_u^Z\longrightarrow K_{Z_b}
 \label{v:eq:span-nat}
\end{equation}
is a strictly commuting comparison of spans whose three vertical maps are
homotopy equivalences. It retains the exact domain identity
\begin{equation}
 D_{vu}^Z=\{\sigma\in D_u^Z:f_u[\sigma]\in D_v^Z\},
 \qquad
 D_{vu}^L=\{\sigma\in D_u^L:f_u[\sigma]\in D_v^L\}.
 \label{v:eq:domain-composition}
\end{equation}
If, additionally, intersection generation and the support flag~\eqref{v:eq:universal-object} hold,
the corresponding concept-lattice actions and supported-domain spans are
identified by extent restriction.
\end{proposition}

\begin{proof}
The full-complex assertions follow from Theorem~\ref{v:thm:three}.
Let $\rho_a$ be its row-dominating retraction. If $e\in E_b$ witnesses
$f_u[\sigma]$, then $g_u(e)$ belongs to every row in $\sigma$.
It therefore belongs to every row of $\rho_a[\sigma]$, so $e$ witnesses
$f_u\rho_a[\sigma]$ as well. The same reasoning applies to
$\sigma\cup\rho_a[\sigma]$. Thus $\rho_a$ restricts to a contiguous
retraction of $D_u^Z$ onto $D_u^L$. The inclusions themselves commute
strictly with $f_u$ because $L$ is invariant; the chosen retractions need
not commute with the operations.

Equation~\eqref{v:eq:domain-composition} is the established Chu exact-domain law: a witness after
$f_vf_u$ pulls back along $g_v$, and therefore already witnesses survival
of the intermediate face. The same argument works on $L$. This is not
an independent new partial-action theorem.

For the final assertion, corresponding extents have identical closed
intents. The intent of the operation applied to a closed extent with
intent $I$ is $g_u^{-1}[I]$ in both contexts. The equality of intents
identifies the full concept actions and the condition that this new
intent be nonempty, which is precisely the supported domain condition.
\end{proof}

The same conclusions are natural for an increasing family $E_t$ of
selected tests when row domination is verified for their union. One
choice of $\rho$ works at every level. If intersection representations
are nonempty, restriction of their row equalities to $E_t$ proves
intersection generation and the support condition at every level as well.
Thus the resulting persistence modules are isomorphic. This does not
supply a canonical homology endomorphism for a partial operation: the
inclusion from its domain need not induce an isomorphism within either
individual span.

\subsection{Consequences for the two legs of a reduced wire}
For designated tests $W$ on $Q_P$, pullback gives relations
\begin{equation}
 pR_Pw\ \Longleftrightarrow\ \beta(p)R_Lw
 \ \Longleftrightarrow\ \kappa\beta(p)R_Qw.
 \label{v:eq:exact-realized}
\end{equation}
Because $\beta$ and $\kappa$ are surjective, both are ordinary
sufficient-state maps for these relations. Hence all structural
comparisons proved in the earlier learning note~\cite{LearningNote} apply on this realized
span: extent lattices are isomorphic, supported cores correspond, and
Dowker complexes are homotopy equivalent. These conclusions do not
require $W$ to separate every external behavioral state. If it does not,
the row quotient for $W$ is correspondingly coarser than $Q_P$.

Suppose now that specified extensions to $\bar P$ obey the required
operational laws. Then Theorem~\ref{v:thm:three} gives a useful additional
conclusion:
\begin{quote}
Exact joint profile lifting, $L=\bar P$, is sufficient but unnecessary
for correctness of the ambient row quotient for the \emph{selected test
family}. The exact
weaker condition is that every ambient response row already occur on
$L$. Under that condition the FCA and Dowker comparisons also hold.
To identify this row quotient with full exterior contextual behavior, the
selected family must separate $Q_P$ and be continuation-complete for the
claimed dynamics.
\end{quote}
This is weaker only for extensional conclusions. It cannot be used to
claim that every compatible reduced profile has an actual native state.
Conversely, matching concept lattices and homotopy types is too weak to
establish even the weaker behavioral claim, as Example~\ref{v:ex:triangle}
shows.

\begin{example}[Discarding a port distinction fabricates geometry]
\label{v:ex:phantom}
Take $X=Y=B=\{0,1\}$ with identity ports, $q_X=q_Y=\id$, and
$q_B:B\to\{*\}$. Then
\[
 P=L=\{00,11\},\qquad \bar P=\{00,01,10,11\}.
\]
The reduced pair $01$ is compatible over $Q_B$ but has no actual compatible
lift. Define three ambient tests with columns
\[
 C_0=\{00\},\qquad C_1=\{11\},\qquad C_{\ne}=\{01,10\}.
\]
These are the tests ``both values are zero'', ``both values are one'', and
``the two values disagree''. Their restrictions to $L$ have columns
$\{00\},\{11\},\varnothing$. With identity dynamics they form a
complete continuation-closed operational example.

The source has two behavioral rows, four closed extents and a Dowker
complex consisting of two points. The ambient context has three rows,
five closed extents and two isolated points plus an edge. Its zeroth
Betti number is three instead of two. The extra component is supported
only by an impossible disagreement test outcome. Thus the state-level
failure of joint lifting has a directly calculable structural consequence;
it is not merely a mismatch of unobserved implementation labels.
\end{example}

\begin{example}[Failure of joint lifting need not harm response geometry]
Keep the same wire, but observe only the first coordinate through the two
tests with columns $\{00,01\}$ and $\{10,11\}$. Every ambient row is
represented on the diagonal. Both behavioral quotients have two states
and both extent lattices have four elements. The realized Dowker complex
has two points; the ambient one has two disjoint edges, and the inclusion
is a homotopy equivalence. Nonetheless neither $01$ nor $10$ is natively
realizable. This is the distinction between row realization and joint
profile lifting, in the same small wire.
\end{example}

\begin{example}[Concepts and nontrivial topology can miss a new behavior]
\label{v:ex:triangle}
Let $L=\{0,1,2\}$ and $E=\{a,b,c\}$ with rows
\[
 r(0)=\{a,c\},\qquad r(1)=\{a,b\},\qquad r(2)=\{b,c\}.
\]
Adjoin a new object $z$ with $r(z)=\{a\}=r(0)\cap r(1)$, and put
$Z=L\cup\{z\}$. The old object complex is the triangle boundary.
The new one replaces the edge $01$ by the filled triangle $01z$, keeping
the other two edges. It deformation retracts to the old circle.
The extent restriction is an isomorphism because the new row is a
nonempty intersection of old rows. Both full concept lattices have eight
elements and both supported cores correspond.

Nevertheless the row quotient grows from three states to four: the new
response profile $\{a\}$ was not attained by any old object. Identity
dynamics make all these statements valid for complete continuation tests.
Thus even agreement of the full concept lattice, supported concept poset,
and a noncontractible Dowker homotopy type does not establish contextual
behavioral adequacy. The marked native state generators matter.
\end{example}

\begin{example}[Why the empty intersection is a real boundary]
Let the two realized rows be $\{a\}$ and $\{b\}$ and adjoin the row
$\{a,b\}$, the empty intersection. The full extent lattices are both
four-element diamonds. The realized complex is two isolated points, while
the new universal object joins them by a two-edge path. Condition~\eqref{v:eq:universal-object}
fails. This is the object-side counterpart of the already known universal
test obstruction, included only to check the exact hypothesis of
Theorem~\ref{v:thm:three}.
\end{example}

\subsection{Preserving a response witness is not preserving its native fibre}
The word ``witness'' has two distinct uses. A Dowker witness is a labelled
test $e$ shared by several states. A source implementation may instead
retain an actual set or category of certificates $\mathcal W(x,e)$ for
the incidence $xRe$. Boolean semantics keeps only
\begin{equation}
 xRe\quad\Longleftrightarrow\quad\mathcal W(x,e)\ne\varnothing.
 \label{v:eq:support}
\end{equation}
The exact sufficient-state comparisons preserve the specified labelled
incidence. The weaker domination comparison preserves only its stated
geometric conclusions. Neither recovers the discarded native fibre.

For an exact comparison of certificate sets along $q:X\twoheadrightarrow
Q$, the relevant additional data are a family $\overline{\mathcal W}$
over $Q\times E$ and a \emph{cartesian} square
\begin{equation}
\begin{array}{ccc}
 \mathcal W&\longrightarrow&\overline{\mathcal W}\\
 \downarrow&&\downarrow\\
 X\times E&\xrightarrow{\ q\times\id_E\ }&Q\times E.
\end{array}
 \label{v:eq:cartesian}
\end{equation}
Equivalently, the specified maps on each fibre are bijections
$\mathcal W(x,e)\cong\overline{\mathcal W}(qx,e)$.
For retained witness categories the appropriate comparison must be stated
at that categorical level, including its morphisms; replacing it by a
bijection of isomorphism classes is an additional forgetting step.

\begin{proposition}[Boolean adequacy does not imply certificate descent]
There is a two-state, one-test example in which the behavioral quotient,
full concept lattice and Dowker homotopy type are all preserved by state
collapse, but no comparison of the form~\eqref{v:eq:cartesian} exists.
\end{proposition}
\begin{proof}
Let $X=\{x_0,x_1\}$, $Q=\{*\}$, and $E=\{e\}$, with the test true
at both states and identity dynamics. Set
$|\mathcal W(x_0,e)|=1$ and $|\mathcal W(x_1,e)|=2$.
The Boolean context has one row class and one concept; its object complex
is an edge and collapses to a point. A cartesian square would require the
one target fibre to be in bijection with both a one-element and a
two-element set, which is impossible.
\end{proof}

Counts are only the simplest obstruction. If native sequential composition
matches an output and an input interface, its certificates are a pullback
$\mathcal W_1\times_B\mathcal W_2$. A certificate-level reduction must
retain the actual map from that pullback to the proposed reduced one;
its invertibility cannot be inferred from separate incidence equivalences.
These are additional lifting requirements beyond incidence equivalence,
not consequences of concept or homotopy adequacy. The structural comparison
keeps them visible when the automata/FCA/Dowker views are attached to a
contextual learner quotient.




\section{What has been generalized and what has not}
The common construction now has an explicit compositional test. A supplied
operational doctrine determines continuation observations. Under the regular
image hypotheses these produce a reflector on state systems. A supplied
wiring diagram then produces the image span of
Theorem~\ref{g:thm:comparison}. Joint lifting and external separation are
independent conditions for that span to be an isomorphism. Actual
regular-epi-preserving wiring functors yield coherent comparisons by
Theorem~\ref{g:thm:coherence}; this is an established minimization mechanism
with its premises made visible.

The native clauses are substantive rather than automatic. Polynomial
machines construct the quotient operations on arbitrary sets and give the
precise backward-request wiring of target-based learners.
Theorem~\ref{s:thm:source-test} gives a source differential certificate
for a regular smooth response state, and
Theorem~\ref{s:thm:learner-wiring} proves that certified reductions compose
under serial and parallel cotangent-feedback learning.
Theorems~\ref{s:thm:gradient} and \ref{s:thm:wiring} add the exact optimizer
condition and its compatibility with objective assembly. The nonlinear Euler
counterexample shows that even this geometric condition does not settle the
actual finite update. The affine class supplies a broad exact finite-step
family rather than leaving the positive result at the level of assumptions.

The structural views are tied to the same maps. On the realized image,
response pullback preserves the selected conceptual and Dowker
characterizations. On the full reduced pullback, equality of row quotients for the selected
family requires realization of every response row; concept adequacy
requires only generation by intersections, and geometric adequacy can be
weaker still. Thus a favorable FCA lattice or persistence calculation is
not a substitute for proving native or behavioral adequacy. A native
certificate fiber is yet another level and needs its actual comparison.

\subsection{Relation to the preceding optimization manuscripts}
The fixed-budget manuscript \cite{BudgetNote} characterizes a specified
native stochastic realization problem by exact zeros of an attained quartic
sum of squares, and proves gradient assembly from local residual factors.
It does not turn an optimizer's positive stationary loss into a feasibility
certificate. Its feasible zeros represent implementations; arbitrary
optimizer states need not be implementations at all. Applying the present
smooth theorem to that encoding would require an actual regular response
map and closure on the states traversed by the optimizer, including off-zero
states. Those hypotheses have not been proved here for that encoding.

The interface-change manuscript \cite{InterfaceNote} separates transport of
feasible realizations, refactorization of a transported tensor, and the larger
saving possible when the target problem admits new tensors. It also
separates zero-set transport, loss equality, gradient transport, and finite
update transport, and proves failures for its actual quartic encoding.
The smooth theorems here do not erase its boundary or no-$C^1$ obstructions.
They explain a general positive class in which the extra identities hold.
The all-loss condition specializes to $DqDq^{\mathsf T}$ descending in the
Euclidean case, already present in the learning-generated characterization
note \cite{LearningNote}; the result is not claimed as a new geometric
principle.

The previous learning note also treats scalar deep linear networks at
arbitrary finite depth, including singular states and exact finite updates.
The present regular smooth theorem applies to arbitrary input/output
dimensions and smooth evaluators when its hypotheses hold, but does not
subsume that singular global classification. The advance in scope is a
conditional doctrine-level bridge, not a claim that a finite example has
solved general realization or language discovery.

\subsection{Contextads and source realization}
Contextads organize context-dependent arrows and their composition
\cite{contextads}. The programme's native constructions additionally retain
actual endpoint squares, occurrence structure, noninvertible cells, and
witnessed factorizations \cite{RelativeNote,RecoveryNote}. A set-valued
behavior quotient cannot reconstruct these by assigning its objects a new
name. To lift the current comparison to that level, one needs the native
interpretation, realizations of the quotient objects, and native comparison
cells satisfying the required unit, composition, and base-change laws.
The full-subcategory criterion above is one sufficient case. It is not
asserted that every native contextad satisfies it.

There is a concrete route in the finite deterministic case: the behavioral
relation on a disjoint union of two machine tables is a computed bisimulation.
When the tables are realized as single-owner deterministic primitives in the
programme's controlled-recovery doctrine, public events and costs agree,
all additional client primitives preserve the relation, and the initial
states are related, the one-point outcome couplings satisfy that doctrine's
local replacement criterion. Its existing theorem then supplies the
admitted adaptive-client and paired-history conclusion \cite{RecoveryNote}.
Those native consequences are inherited from that theorem under its actual
premises; they do not follow from the coalgebra quotient alone. Unrepresented
shared state, changed authority, or possibility-only branches are not
licensed by this instantiation.

\subsection{Closest prior results and the claim boundary}
For composition and minimization, \cite[Propositions 14, 43, 46 and
Corollary 47]{RSW} is a direct antecedent, not merely related background.
For the polynomial losslessness criterion, \cite[Theorem 3.9]{bartels}
and \cite[Lemma 4.6]{bonchi} already supply the key injective-transformation
principle; the two/three-state converse is an elementary specialization,
with degeneracy warnings already present in the former paper.
Factorization-system minimization is treated generally in \cite{wissmann}.
Turi--Plotkin's distributive-law framework supplies a stronger established
syntax/behavior compatibility theory when its hypotheses are met
\cite{turi}; invoking its name would not prove those hypotheses for a new
wiring language. Fong--Sarazola likewise give a general recipe for
compositional black-box functors \cite{blackbox}. Black boxing does not
by itself establish minimal native realization.

For learners, \cite{FST,CGGWZ,spivak} already provide compositional
interfaces and state-map preservation laws. Smooth projectability and
exact lumpability have direct predecessors in \cite{HA}; intrinsic learning
dynamics, metrics, and initialization-dependent invariants in \cite{MPG}.
The generalized gradient hypergraph functor \cite{HFK} transports objective
covectors and assembles fields contravariantly. Its field reindexing is
$v\mapsto\varphi^\dagger v\varphi$, not the covariant state projection
$dq\,v$. The categorical reduction functor in
Theorem~\ref{s:thm:functor} has deliberately restricted arrows satisfying
the extra cometric law. Neither theorem proves global optimization,
canonical choice of minimizers, or commutation of arbitrary nonlinear
Euler discretizations.

The result is therefore a general, proved framework for organizing and
certifying contextual reduction across wiring, native learning, and
structural characterizations.
This develops compositional use of characterizations. Formation and revision
of the operational language, and stabilization of scientific or mathematical
theories, remain broader questions. Historical novelty of a more specific
synthesis has not been established. The next substantive theorem should
verify or generate the regular response coordinates, boundary-covering
maps, and external separating protocols in an independently specified
native learner grammar, or prove why they cannot exist. An approximate
version would additionally need explicit error, horizon, and resource
bounds. These are precise remaining mathematical problems, not conclusions
obtained by declaring all contexts already available.


\begin{thebibliography}{99}
\bibitem{RSW} R. Rosebrugh, N. Sabadini, and R. F. C. Walters.
\emph{Minimal realization in bicategories of automata}.
Mathematical Structures in Computer Science 8(2) (1998), 93--116.
In particular Propositions 14, 43, 46 and Corollary 47.
\url{https://doi.org/10.1017/S0960129597002454}.
\url{https://www.mta.ca/~rrosebru/articles/mnrl.pdf}.

\bibitem{BudgetNote}
\emph{Fixed-Budget Native Realization and Compositional Optimization}.
Prior project note, November 2026. The native occurrence and statewise
continuation constraints, exact fixed-budget encoding, and separation of
gradient assembly from optimization guarantees are retained here.

\bibitem{InterfaceNote}
\emph{Interface Change and Realization Optimization}.
Prior project note, November 2026. In particular the separation of
feasible transport, loss transport, and actual field/update transport.

\bibitem{rutten} J. J. M. M. Rutten.
\emph{Universal coalgebra: a theory of systems}.
Theoretical Computer Science 249 (2000), 3--80.
\url{https://ir.cwi.nl/pub/48}.

\bibitem{spivak} D. I. Spivak. \emph{Learners' Languages}.
EPTCS 372 (2022), 14--28; revised arXiv version 3, June 2025.
\url{https://arxiv.org/html/2103.01189}.

\bibitem{calf} G. van Heerdt, M. Sammartino, A. Silva.
\emph{CALF: Categorical Automata Learning Framework}.
CSL 2017, LIPIcs 82, article 29.
\url{https://drops.dagstuhl.de/entities/document/10.4230/LIPIcs.CSL.2017.29}.

\bibitem{bartels} F. Bartels, A. Sokolova, and E. de Vink.
\emph{A hierarchy of probabilistic system types}.
Electronic Notes in Theoretical Computer Science 82(1) (2003).
Theorem 3.9 and Section 3 counterexample discussion.
\url{https://www.cs.uni-salzburg.at/~anas/assets/papers/CMCS03-Hierarchy.pdf}.

\bibitem{bonchi} F. Bonchi, A. Silva, A. Sokolova.
\emph{Distribution Bisimilarity via The Power of Convex Algebras}.
Logical Methods in Computer Science 17(3), article 10 (2021), Lemma 4.6.
\url{https://doi.org/10.46298/LMCS-17(3:10)2021}.

\bibitem{contextads} M. Capucci, D. J. Myers.
\emph{Contextads as Wreaths; Kleisli, Para, and Span Constructions as Wreath Products}.
arXiv:2410.21889 (2024).
\url{https://arxiv.org/abs/2410.21889}.

\bibitem{RelativeNote} Author-supplied manuscript.
\emph{One relative construction for context content and effects},
1 November 2026, main theorem and realization construction.

\bibitem{RecoveryNote} Author-supplied manuscript.
\emph{Contextual Refinement and Controlled Recovery},
1 November 2026, Sections 2--5 and the native paired-history construction.

\bibitem{FST}
Brendan Fong, David I. Spivak, and R\'emy Tuy\'eras.
\emph{Backprop as Functor: A compositional perspective on supervised
learning}. LICS 2019. In particular Sections II and VII-C.
\url{https://arxiv.org/html/1711.10455v3}.

\bibitem{CGGWZ}
G. S. H. Cruttwell, Bruno Gavranovi\'c, Neil Ghani, Paul Wilson, and Fabio
Zanasi. \emph{Categorical Foundations of Gradient-Based Learning}.
ESOP 2022; arXiv version 2, 2021.
\url{https://arxiv.org/abs/2103.01931}.

\bibitem{HA}
Leonhard Horstmeyer and Fatihcan M. Atay.
\emph{Characterization of exact lumpability for vector fields on smooth
manifolds}. Differential Geometry and its Applications 48 (2016), 46--60.
\url{https://arxiv.org/html/1607.01237v2}.
\url{https://doi.org/10.1016/j.difgeo.2016.06.001}.

\bibitem{HFK}
Tyler Hanks, James Fairbanks, and Matthew Klawonn.
\emph{Generalized Gradient Descent is a Hypergraph Functor}.
Electronic Proceedings in Theoretical Computer Science 429 (2025), 217--233.
Theorems 3.4 and 3.5, Corollary 3.5.1.
\url{https://arxiv.org/html/2403.19845v2}.
\url{https://doi.org/10.4204/EPTCS.429.12}.

\bibitem{MPG}
Sibylle Marcotte, Gabriel Peyr\'e, and R\'emi Gribonval.
\emph{Intrinsic training dynamics of deep neural networks}.
arXiv:2508.07370, version 2. In particular Definitions 2.10 and 2.15,
Lemma 2.11, and Theorem 2.14.
\url{https://arxiv.org/html/2508.07370v2}.

\bibitem{LearningNote}
\emph{Learning generated characterizations}.
Earlier project manuscript dated 5 November 2026.
Source compared: \nolinkurl{learning_generated_characterizations.tex}.

\bibitem{gp} J. Gratus and T. Porter,
\emph{A geometry of information, I: Nerves, posets and differential forms},
Dagstuhl Seminar Proceedings 04351 (2005), Proposition 1.
\href{https://drops.dagstuhl.de/storage/16dagstuhl-seminar-proceedings/dsp-vol04351/DagSemProc.04351.6/DagSemProc.04351.6.pdf}{Primary text}.

\bibitem{ayzenberg} A. Ayzenberg,
\emph{Topology of nerves and formal concepts} (2019), Proposition 5.7.
\href{https://arxiv.org/abs/1911.05491}{arXiv:1911.05491}.

\bibitem{brun} M. Brun and L. M. Salbu,
\emph{The rectangle complex of a relation}, Mediterranean Journal of
Mathematics 20 (2023), article 7.
\href{https://doi.org/10.1007/s00009-022-02213-0}{Published article}.

\bibitem{CommonNote} Earlier project manuscript.
\emph{One operational construction for residual automata, formal concepts,
and topology}. 4 November 2026. Source supplied as
\texttt{common\_generation\_note.tex}.

\bibitem{wissmann} T. Wi\ss mann.
\emph{Minimality notions via factorization systems and examples}.
Logical Methods in Computer Science 18(3:31) (2022).
\url{https://lmcs.episciences.org/10014}.
\url{https://arxiv.org/abs/2202.05701}.

\bibitem{turi} D. Turi and G. Plotkin.
\emph{Towards a mathematical operational semantics}. LICS 1997.
\url{https://homepages.inf.ed.ac.uk/gdp/publications/Math_Op_Sem.pdf}.

\bibitem{blackbox} B. Fong and M. Sarazola.
\emph{A recipe for black box functors}. Theory and Applications of
Categories (2020). Theorems 3.21 and 5.2.
\url{https://arxiv.org/abs/1812.03601}.
\end{thebibliography}
\end{document}
